\chapter{NNLL resummation in $Z$+jet events with the ARES method}
\label{ch:hhZjet}

\section{Introduction}
\label{sec:intro}

To date, no general framework has been established for performing
next-to-next-to-leading logarithmic (NNLL) resummation of jet
observables in hadronic collisions. Considerable effort is currently 
being devoted to extending parton showers
to this accuracy, but the current state of the art remains at NLL,
demonstrated by \textsc{Sherpa}~\cite{Herren:2022jej,Sherpa:2024mfk}
and \textsc{PanScales}~\cite{vanBeekveld:2022zhl,vanBeekveld:2022ukn,
vanBeekveld:2023chs}. Within Soft-Collinear Effective Theory (SCET),
resummation has been achieved up to N$^3$LL accuracy for a number of
observables in hadronic collisions~\cite{Becher:2008cf, Abbate:2010xh, Alioli:2023rxx}. However, these calculations rely on
observable-specific factorisation theorems, limiting their applicability
to broader classes of observables. Although recent progress has been made
towards a more general framework~\cite{vanBeekveld:2025zjh}, no general formalism applicable to a
wide range of event shapes currently exists.

In this chapter, we extend the \textsc{ARES} formalism for
NNLL resummation to hadronic collisions, considering the process $pp\to Z+\text{jet}$, 
where the $Z$ boson recoils against a jet with high transverse momentum. We
find that the structure of the NNLL resummation is formally identical
to that for three-jet events in $e^+e^-$ annihilation~\cite{Arpino:2019ozn}, with the same
master formula but different ingredients. The main new feature is the
treatment of initial-state radiation, whose NNLL contributions involve
convolutions of splitting functions with parton distribution functions.
These terms have no analogue in the $e^+e^-$ case and constitute the key
extension required for applications to hadronic collisions.

As a first application of the formalism, we consider the event-shape
observable one-jettiness. This observable has been extensively studied
and resummed to N$^3$LL accuracy~\cite{Alioli:2023rxx}, providing a
useful benchmark against which our results can be validated.
Furthermore, one-jettiness is an additive observable
(see Sec.~\ref{subsec:additive}), making it particularly straightforward
to calculate within \textsc{ARES}. The aim of this
application is therefore not to provide new phenomenological
predictions, but rather to demonstrate the implementation of the
extension to hadronic collisions and validate the resulting resummation
against known results for one-jettiness~\cite{Gaunt:2015pea}.
The structure of the initial-state radiation contributions is validated independently by comparing
the corresponding NNLL terms with the results obtained for Lund tree shapes
in Ref.~\cite{vanBeekveld:2025zjh}.

Unlike the NLL discussion presented in the previous chapter, we do not
derive the NNLL master formula from first principles, as this is not the
primary focus of this work. Instead, we focus on its extension to
hadronic collisions, in particular the treatment of initial-state
radiation. The derivation of the general NNLL master formula can be
found in Refs.~\cite{Banfi:2014sua,Banfi:2018mcq,Arpino:2019ozn}, to
which we refer the reader for a detailed discussion.

The remainder of this chapter is organised as follows. We first present
the \textsc{ARES} NNLL master formula for hadronic collisions and derive
the corresponding initial-state radiation contributions. We then apply
the formalism to one-jettiness and conclude with a validation of the
results against existing calculations in the literature.

\section{General setup}
\label{sec:setup}

We consider the production of a $Z$ boson and a hard jet in hadron collisions,
\begin{equation}
  \label{eq:Z+jet}
  h_1(P_1)+h_2(P_2) \to Z(q) + \mathrm{jet}(p_J)+\dots\,,
\end{equation}
where $h_1$ and $h_2$ denote the incoming hadrons with momenta $P_1$ and $P_2$, and the dots represent accompanying hadrons. At Born level, the process proceeds through two incoming partons with momenta $p_1=x_1P_1$ and $p_2=x_2P_2$, producing a $Z$ boson with momentum $q$ and a single hard parton with momentum $p_3$. 
Beyond Born level, additional QCD radiation causes the hard parton to recoil, modifying its momentum to $\tilde p_3$. The recoiling parton is subsequently clustered together with the accompanying radiation to form the observed jet with momentum $p_J$. In the Born limit, the two coincide, so that $p_J=p_3$.

We consider event shapes satisfying the conditions set out in
Sec.~\ref{sec:obs-req}, denoted generally as functions $V(q_1,\dots,q_n)$ of the
hadron momenta $q_1,\dots,q_n$.
As an example, we consider
one-jettiness, an observable that has been widely studied in SCET.  It corresponds to the 
$N=1$ case of the more general $N$-jettiness observable~\cite{Stewart:2010tn}, 
which quantifies how well the final state can be described by $N$ narrow jets. It is
defined as
\begin{equation}
 \label{eq:T1}
 \mathcal{T}_1\equiv \sum_{i} \min\left\{\frac{2 \bar{p}_1\cdot q_i}{Q_1}, \frac{2 \bar{p}_2\cdot q_i}{Q_2}, \frac{2 p_J \cdot q_i}{Q_3}\right\}\,,
\end{equation}
where $q_i$ is the four-momentum of a \emph{hadron}, and
$Q_1, Q_2, Q_3$ are three momentum scales. The two ``incoming'' momenta, 
$\bar p_1$ and $\bar p_2$ are defined as 
\ABquestion{Still don't quite understand the difference between $\bar{p}_i$ and $p_i$
and what we mean by the quotation marks around ``incoming''  are meant to indicate?
I would not know what to say what these objects are if asked.}
\begin{equation}
	\bar{p}_1^\mu=\frac{1}{2}n_2\cdot (p_J+q)n_1^{\mu}\,,\quad\quad
	\bar{p}_2^\mu=\frac{1}{2}n_1\cdot (p_J+q)n_2^{\mu}\,,
\end{equation}
where
\begin{equation}
  \label{eq:ni-incoming}
  n_1 = (1,0,0,1)\,,\qquad n_2 = (1,0,0,-1)\,.
\end{equation}
When $\mathcal{T}_1\to 0$, we are in the \emph{one-jet region}, in which
the outgoing jet $p_J$ recoils against the $Z$ boson. In this limit,
$\bar p_1$ and $\bar p_2$ coincide with the Born momenta $p_1$ and
$p_2$.

We also introduce the normalised observable
$\tau_1\equiv \mathcal{T}_1/M$, where $M$ denotes a generic hard scale. 
Choosing
\begin{equation}
  \label{eq:Qi-inv}
  Q_1 = \frac{2 \bar{p}_1\cdot p_J}{p_{T,J}}\,,  \quad Q_2 = \frac{2 \bar{p}_2\cdot p_J}{p_{T,J}}\,,
  \quad Q_3 = 2 p_{T,J}\,,
\end{equation}
where $p_{T,J}$ is the transverse momentum of the outgoing jet with
respect to the beam, provides a definition of $\mathcal{T}_1$ that is
manifestly invariant under boosts along the beam direction,
\begin{equation}
	\label{eq:T1-binv}
	\mathcal{T}_1\equiv  \sum_{i} k_{t,i}\min\{e^{(\eta_i-\eta_J)}, e^{-(\eta_i-\eta_J)}, \cosh(\eta_{i}-\eta_J)-\cos(\phi_{i}-\phi_J)\}\,,
\end{equation}
where $k_{t,i}$ and $\eta_i$ denote the transverse momentum and rapidity
of hadron $q_i$ with respect to the beam, while $\eta_J$ is the rapidity of the jet. 
To reduce the impact of experimental systematics associated with the measurements of particle momenta, 
we can set $M=H_T$, with $H_T$ the scalar sum of all hadron transverse
momenta. This leads to the dimensionless observable
\begin{equation}
	\label{eq:tau1-binv}
	\tau_1\equiv  \sum_{i} \frac{k_{t,i}}{H_T}\min\{e^{(\eta_i-\eta_J)}, e^{-(\eta_i-\eta_J)}, \cosh(\eta_{i}-\eta_J)-\cos(\phi_{i}-\phi_J)\}\,.
\end{equation}

As discussed in Sec.~\ref{sec:problem-spec}, we consider the cumulant of
an event shape, obtained by imposing the constraint
$V(q_1,q_2,\dots,q_n)<v$. For one-jettiness, the one-jet region corresponds
to the limit $v\to0$. To fix the Born kinematics in this limit, we
consider the cross section differential in $\eta_J$, and in the transverse momentum $q_T$ 
and rapidity $y_Z$ of the $Z$ boson. The differential cross section then takes the factorised form
of Eq.~\eqref{eq:Sigma-H-fact}, specialised to the present process, \ABquestion{Can we say it takes the factorised form of Eq.~\eqref{eq:Sigma-H-fact}}
\begin{equation}
  \label{eq:dsigma-def}
  \frac{d\sigma(v)}{dq_T dy_Zd\eta_J} \simeq \sum_{i,j} \int_0^1 \!\! dx_1\, f_i(x_1,\mu_F)\int_0^1 \!\! dx_2\, f_j(x_2,\mu_F)\frac{d\hat \sigma_{ij}}{dq_Tdy_Zd\eta_J} \Sigma_{ij}(\{p_1,p_2,p_3\},v)\,,
\end{equation}
In the above equation, we have introduced the partonic cross section, $d\hat{\sigma}_{ij}$,
for the production of a $Z$ boson and an extra jet via the partonic
subprocess $ij$, where $i$ represents a quark, an antiquark of any
light flavour ($u,d,c,s,b$, considered massless), or a gluon. 
The expression in Eq.~\eqref{eq:dsigma-def}
requires the construction of a mapping between the full event -- made up
of the three hard partons $\{\tilde p\}\equiv\{\tilde p_1,\tilde p_2,\tilde p_3\}$ and an arbitrary number of soft and/or collinear emissions $\{k_i\}$ -- and
a Born event with only three hard partons $\{p_1,p_2,p_3\}$.  \ABquestion{Is it necessary to expand on
what this mapping is? Can I at least say the mapping is in the 3-jet paper and cite that or something? I don't understand
what this mapping is here.}
The difference between the approximation on the right-hand side of
Eq.~\eqref{eq:dsigma-def} and the full expression for
$ d\sigma(v)/(dq_T dy_Z d\eta_J)$ is a well-defined expansion in powers of the
strong coupling $\alpha_s$, with coefficients vanishing as a power of
$v$, potentially enhanced by logarithms of $v$.

The cumulative distribution $\Sigma_{ij}(\{p_1,p_2,p_3\},v)$ contains
logarithms of $v$ that become large in the one-jet region. The
largest of these contributions are double logarithms of the form
$(\alpha_s L^2)^n$, with $L\equiv \ln(1/v)$. Neglecting 
coherence-violating effects, if $V(\{\tilde p\},\{k_i\})$ is rIRC
safe, for each Born phase space point, in the region
$\alpha_s L\sim 1$, $\Sigma_{ij}(\{p_1,p_2,p_3\},v)$ can be
reorganised as follows
\begin{equation}
  \label{eq:log-counting}
  \Sigma_{ij}(\{p_1,p_2,p_3\},v) = e^{L g_1(\alpha_s L)}\left[G_2(\alpha_s L)+ \alpha_s G_3(\alpha_s L)+\dots\right]\,.
\end{equation}
As illustrated in Table.~\ref{tab:log-accuracy} the function $g_1(\alpha_s L)$ resums all leading logarithmic (LL)
terms, $\alpha_s^n L^{n+1}$, in $\ln\Sigma_{ij}$. For rIRC safe
observables, $g_1$ exponentiates. The function $G_2(\alpha_s L)$
resums next-to-leading logarithmic (NLL) contributions
($\alpha_s^n L^n$) to $\ln\Sigma_{ij}$, and so on. In this work, we
aim at NNLL accuracy, i.e.\ calculating $g_1$, $G_2$ and $G_3$. 


\section{General form of NNLL resummation}
\label{sec:NNLL-general}

The main differences between event shape resummation in hadronic
collisions and $e^+e^-$ annihilation are
\begin{itemize}
\item the presence of initial-state radiation;
\item the fact that there is no single natural hard scale, $Q$. \ABquestion{This is confusing
as $Q$ is also the resummation scale? Should we just remove the $Q$ here?}
\end{itemize}
Nevertheless, as will be made clear in Sec.~\ref{sec:NNLL-initial}, the structure of the NNLL resummation for $\Sigma_{ij}(\{p_1,p_2,p_3\},v)$ is the same as for $e^+e^-$ annihilation, although several ingredients appearing in the
resummation formula are modified to account for initial-state radiation.
 In this section, we present the general structure of the NNLL resummation, and discuss the physical meaning of each ingredient. 
 For each partonic channel $ij$, we have 
\begin{multline}
  \label{eq:Sigma-NNLL}
  \Sigma_{ij}(\{p_1,p_2,p_3\},v)= e^{-R_{\rm NNLL}(v)} \frac{f_i(x_1,v^{\frac{1}{a+b_1}}Q)}{f_i(x_1,\mu_F)} \frac{f_j(x_2,v^{\frac{1}{a+b_2}}Q)}{f_j(x_2,\mu_F)}
  \times \\ \times
  \left(\mathcal{F}^{(ij)}_{\rm NLL}(\lambda)\left(1+\frac{\alpha_s(Q)}{2\pi}H_{ij}^{(1)}(Q,\mu_R)
      +\frac{\alpha_s(Q v^{\frac{1}{a+b_\ell}})}{2\pi}\sum_{\ell}C_{\mathrm{hc},\ell}^{(1)}\right) +\frac{\alpha_s(Q)}{\pi} \delta\mathcal{F}^{(ij)}_{\rm NNLL}(\lambda)\right)\,,
\end{multline}
where $\mu_R,\mu_F$ are the renormalisation and factorisation
scales in the $\overline{\rm MS}$ scheme, and $\beta_0$ and $\lambda$
are defined in Eqs.~\eqref{eq:beta0} and~\eqref{eq:lambda} respectively.
The parameters $a_\ell, b_\ell$ are determined by the general behaviour 
of an IRC safe observable $V$ in the presence of a single soft emission collinear to leg $\ell$,
as given in Eq.~\eqref{eq:general-sc-V}. For hadronic collisions, the scale $Q$ is an arbitrary resummation
scale.

The physical meaning of all ingredients in
Eq.~\eqref{eq:Sigma-NNLL} can be recovered from
Refs.~\cite{Arpino:2019ozn,Banfi:2018mcq,Banfi:2004yd}. We give
below a short description of each, highlighting the differences with
respect to the $e^+e^-$ case:

\begin{itemize}
\item The NNLL radiator $R_{\rm NNLL}(v)$ includes all virtual
  corrections of \emph{soft and/or collinear origin}. These are divergent,
  but their divergences cancel with those of real radiation. In
  order to cancel the singularities of virtual corrections analytically, the
  radiator also includes a subset of real corrections for suitably
  defined unresolved emissions. The procedure to define those
  emissions is not unique (cf.\ Sec.\ref{sec:NLL-resum}). A proposal adequate for three emitting legs is given in
  Ref.~\cite{Arpino:2019ozn}, leading to a master formula for
  $R_{\rm NNLL}(v)$ whose explicit expression depends only on the
  single-emission parameters $a$, $b_\ell$, $d_\ell$, and
  $g_\ell(\phi)$. For completeness, we illustrate how
  the NNLL radiator extends from the NLL result in
  Sec.~\ref{sec:radiator}. It is worth noting that  
  with the knowledge we have acquired so far, the radiator 
  could be computed for a generic observable at N$^3$LL accuracy. 

  \item Parton distribution functions enter through the ratios
  $f_i(x_\ell,v^{\frac{1}{a+b_\ell}}Q)/f_i(x_\ell,\mu_F)$,
  where $f_i(x_\ell,\mu_F)$ is the parton distribution function
  corresponding to the flavour of the parton $p_\ell$ entering the hard scattering, with
  momentum fraction $x_\ell$. Between the factorisation scale $\mu_F$ and the veto scale
  $v^{1/(a+b_\ell)}Q$, the incoming parton can emit \emph{hard-collinear}
  gluons that are too collinear to be resolved by the observable.
  These unresolved emissions are absent from
  $\mathcal{F}_{\rm NLL}$ and instead they modify the effective PDF seen
  by the hard process. The ratio $f_i(x_\ell,v^{1/(a+b_\ell)}Q)/
  f_i(x_\ell,\mu_F)$ resums precisely this effect via DGLAP
  evolution from $\mu_F$ down to the veto scale. This contribution is absent in
  $e^+e^-$ annihilation, where there are no incoming partons, and
  is a genuinely new ingredient in hadronic collisions.

\item Resolved \emph{soft and collinear} emissions widely separated in angle,
  together with the appropriate virtual corrections, build up the
  NLL function $\mathcal{F}_{\rm NLL}(\lambda)$, derived in
  Sec.~\ref{sec:NLL-resum} and given in Eq.~\eqref{eq:FNLL-formula}.
  This ingredient is unchanged from the $e^+e^-$ case.

\item Finite virtual corrections of relative order $\alpha_s$ are
  encoded in the hard coefficient $H_{ij}^{(1)}$. This is a term of order
  $\alpha_s$ multiplying the NLL function, $\mathcal{F}^{(ij)}_{\rm NLL}(\lambda)$, 
  hence NNLL. For a generic process, $H_{ij}^{(1)}$ depends on $Q$, kinematic invariants and $\mu_R$. This contribution is process dependent and so will be entirely different between the $\ee$ and hadronic case. The general formula for $H_{ij}^{(1)}$ is derived in Sec.~\ref{sec:H1}; 
   
\item A single unresolved \emph{hard-collinear splitting} of leg
  $\ell$ combines with the corresponding virtual corrections and, in
  the case of initial-state radiation, to the collinear counterterm
  needed to renormalise the parton distribution functions. This gives
  the coefficient $C_{\mathrm{hc},\ell}^{(1)}$, which gives rise to a NNLL
  contribution since it is a term of order $\alpha_s$
  multiplying the NLL function, $\mathcal{F}^{(ij)}_{\rm NLL}(\lambda)$.
  \ABquestion{Can I say: $C_{\mathrm{hc},\ell}^{(1)}$ is the finite piece left over once the
  collinear $1/\varepsilon$ singularities cancel. In $e^+e^-$ annihilation
  this cancellation happens entirely between the real hard-collinear
  splitting and its virtual correction. In hadronic collisions, radiation
  collinear to an incoming leg carries an extra singularity, already
  absorbed into the PDF, that real and virtual corrections alone cannot
  remove. The collinear counterterm removes it
  instead. Combining real, virtual, and, for ISR, the counterterm leaves
  the finite coefficient $C_{\mathrm{hc},\ell}^{(1)}$.}
The final-state contribution is identical to that derived in Ref.~\cite{Arpino:2019ozn}, 
while the initial-state contribution is a new result of this work, derived
 in Sec.~\ref{sec:NNLL-initial}.  
  
  \item All remaining resolved emissions and their associated virtual
  corrections are collected in the NNLL function
  $\delta\mathcal{F}_{\rm NNLL}(\lambda)$. These consist of one
  special \emph{soft and/or collinear} emission accompanied by an
  arbitrary ensemble of soft-collinear emissions widely
  separated in angle, together with the corresponding virtual
  corrections:
   \begin{equation}
    \label{eq:dFNNLL}
    \delta\mathcal{F}_{\rm NNLL} =
    \delta\mathcal{F}_{\rm sc}
    +\delta\mathcal{F}_{\rm hc}
    +\delta\mathcal{F}_{\rm rec}
    +\Delta\mathcal{F}_{\rm rec}
    +\delta\mathcal{F}_{\rm wa}
    +\Delta\mathcal{F}_{\rm wa}
    +\delta\mathcal{F}_{\rm correl}\,.
  \end{equation}
  Each term reflects the region of phase space spanned by the special emission
  and can be expressed as an integral over the soft-collinear measure $\dZ$ with the appropriate weight function $G(\{k_i\})$ 
  (see Eq.~\eqref{eq:dZ-def}). 

The contributions of soft origin ---
$\delta\mathcal{F}_{\rm sc}$, $\delta\mathcal{F}_{\rm wa}$,
$\Delta\mathcal{F}_{\rm wa}$, and $\delta\mathcal{F}_{\rm correl}$
--- are identical to those derived in Ref.~\cite{Arpino:2019ozn} and
receive no additional contributions from initial-state radiation. This follows directly
from the soft matrix-element factorisation described in
Sec.~\ref{subsec:soft-fact}, according to which each soft emission
depends only on the colour charge and direction of the emitting leg,
and is independent of its momentum fraction, flavour, or whether the
emitter is in the initial or final state. Consequently, the soft
corrections are universal for ISR and FSR.

In contrast, the collinear contributions differ between initial- and
final-state radiation. \ABquestion{Is the following correct: 
Emissions collinear to an incoming parton carry
a finite fraction of its longitudinal momentum, giving rise to
collinear singularities that are absorbed into the PDFs and regulated 
through the PDF counterterm.} As a result, the corresponding factorisation involves
splitting functions convoluted with PDFs, and depends explicitly on the
identity and momentum fraction of the incoming parton. This behaviour,
introduced in Sec.~\ref{subsec:collinear-fact}, is reflected in the
hard-collinear contributions --- $\delta\mathcal{F}_{\rm hc}$, $\delta\mathcal{F}_{\rm rec}$, and
$\Delta\mathcal{F}_{\rm rec}$. While these terms retain the same formal
structure as in the $e^+e^-$ case, their ingredients are modified for
ISR through the appearance of PDF convolutions. For the one-jettiness
observable considered in this chapter, the spin-correlation contribution
to $\Delta\mathcal{F}_{\rm rec}$ vanishes for both initial- and
final-state radiation. Therefore, the ISR extension considered here
only requires the derivation of the contributions to
$\delta\mathcal{F}_{\rm hc}$ and $\delta\mathcal{F}_{\rm rec}$.
However, a complete NNLL formulation for generic observables requires
the inclusion of spin-correlation effects. For completeness, the
derivation of the corresponding initial-state contribution is presented in
Appendix~\ref{app:spin-correl}. The remaining initial-state radiation
corrections relevant for the present calculation are derived in
Sec.~\ref{sec:NNLL-initial}.
  
  \end{itemize} 
  
  In the following two subsections we discuss $H_{ij}^{(1)}(Q,\mu_R)$
and $R_{\rm NNLL}(v)$ in more detail. The radiator can in principle
be taken directly from Ref.~\cite{Arpino:2019ozn}, since the Born-level process
$hh\to Z+\text{jet}$ has the same three-leg colour structure as the
$e^+e^-\to q\bar q g$ case. Nevertheless we present its derivation here to
illustrate how the NLL radiator is extended to NNLL and to
establish the conventions used throughout this chapter. The hard
function $H_{ij}^{(1)}$, by contrast, is process specific
and cannot be taken from Ref.~\cite{Arpino:2019ozn}. Its general
formula is derived in the following.


\subsection{The hard coefficient $\mathbf{H_{ij}^{(1)}(Q,\mu_R)}$}
\label{sec:H1}

The hard coefficient $H_{ij}^{(1)}(Q,\mu_R)$ encodes the finite contribution
from one-loop virtual corrections to the hard scattering amplitude.
Its derivation allow us to set the conventions used throughout this chapter and introduce
the web function formalism that also underlies the derivation of the NNLL radiator
(cf.\ Sec.~\ref{sec:radiator}). We follow the notation of Ref.~\cite{Catani:1998bh},
 slightly adapted to our needs.
 
 The starting point for the computation of $H_{ij}^{(1)}$ is the
following representation of UV-renormalised virtual corrections
~\cite{Banfi:2018mcq,Arpino:2019ozn}:
\begin{multline}
  \label{eq:HQ}
  \mathcal{V}(Q;\varepsilon) = H_{ij}(\alpha_s(\mu_R),Q) \exp\left\{\int^{Q^2}\!\!\frac{dk^2}{k^2}\sum_{\ell=1}^3 \gamma_\ell\left(\alpha_s(k;\varepsilon)\right)\right\}\times \\ \times \exp\left\{-\int \frac{d^dk}{(2\pi)^d}\sum_{\ell=1}^3 \sum_{\bar \ell\ne \ell} w_{\ell \bar \ell}(m^2,m^2+\kappa_{\ell\bar \ell}^2;\varepsilon) \Theta\left(\frac 12\ln\frac{Q^2_{\ell\bar \ell}}{m^2+\kappa_{\ell\bar \ell}^2}-|\eta^{(\ell\bar \ell)}| \right) \Theta(Q_{\ell \bar \ell}-\kappa_{\ell\bar \ell})\right\}\,,
\end{multline}
where $Q_{\ell\bar\ell}^2=2p_\ell\cdot p_{\bar\ell}$ and dimensional
regularisation is used to regulate soft and collinear divergences,
with $d=4-2\varepsilon$.  

The virtual corrections are separated into hard-collinear and soft
contributions. The hard-collinear singularities are described by the
$\gamma_\ell$ exponential, where $\gamma_\ell$ is the collinear anomalous
dimension associated with external leg $\ell$. The second exponential
contains the web function $w_{\ell\bar\ell}$, which describes soft
virtual corrections arising from the exchange of an arbitrary number of soft partons between
the dipole $(\ell\bar\ell)$. \ABquestion{I wanted to explain what this $m$ is, can I say:
The origin of the constraints given by the $\Theta$-functions 
are explained in Sec.~\ref{subsubsec:Rs}, where there is an additional $m$-dependence
which is a way of parameterising the multiple emissions?}
The hard function $H_{ij}(\alpha_s,Q)$ contains the finite contribution
remaining after these IR-divergent terms have been extracted. It
therefore represents the process-dependent finite part of the virtual
corrections.

The key ingredients in this representation are:
\begin{itemize}
\item
  the ``hard'' function $H_{ij}(\alpha_s,Q)$ for incoming channel $ij$, which has the following expansion in powers of $\alpha_s(Q)$:
  \begin{equation}
    \label{eq:HQ-exp}
    H_{ij}(\alpha_s,Q) = 1+\frac{\alpha_s}{2\pi}H_{ij}^{(1)}(Q) +\left(\frac{\alpha_s}{2\pi}\right)^2H_{ij}^{(2)}(Q)+\dots 
  \end{equation}
  
\item the $d$-dimensional strong coupling $\alpha_s(k^2;\varepsilon)$, defined as the solution of the renormalisation-group equation
\begin{equation}
  \label{eq:RG-coupling}
  k^2\frac{d\alpha_s}{dk^2} = - \varepsilon \alpha_s+\beta(\alpha_s)\,,
\end{equation}
where $\beta(\alpha_s)$ is the QCD beta function, and renormalisation
is performed in the $\overline{\rm MS}$ scheme. The collinear anomalous dimension $\gamma_\ell(\alpha_s)$ has the expansion 
\begin{equation}
  \label{eq:gammal}
  \gamma_\ell (\alpha_s) = \gamma^{(0)}_\ell \frac{\alpha_s}{2\pi} + \gamma^{(1)}_\ell \left(\frac{\alpha_s}{2\pi} \right)^2+\dots
\end{equation}
where $\gamma_\ell^{(0)}$ is given in Eq.~\eqref{eq:gammal-flav}, and 
 \begin{equation}
  \label{eq:gammal1-flav}
  \gamma_\ell^{(1)} = \left\{
    \begin{split}
      &\frac{C_F}{2}\left(C_F\left(\frac{3}{4}-\pi^2+12\zeta_3\right)+C_A\left(\frac{17}{12}+\frac{11\pi^2}{9}-6\zeta_3\right)-n_f\left(\frac{1}{6}+\frac{2\pi^2}{9}\right)\right)\,,\quad \text{if $\ell$ is a quark,} \\
      & C_A^2\left(\frac{8}{3}+3\zeta_3\right)-\frac{n_f}{2}C_F-\frac{2}{3}n_fC_A \,,\qquad\qquad\qquad\qquad\qquad\qquad\qquad\qquad\qquad\text{if $\ell$ is a gluon}\,;
    \end{split}
    \right.
\end{equation}

\item the ``web'' function
  $w_{\ell \bar \ell}(m^2,m^2+\kappa_{\ell\bar \ell}^2;\varepsilon)$
  which consistently resum the contribution of multiple soft virtual
  corrections from the dipole $(\ell\bar\ell)$. Physically, the web function is the weight associated
  with the exchange of a single gluon of mass $m$, inclusive with respect to
  all its splittings. Since Feynman rules are invariant under
  rescaling of the emitters' four-momenta, the web function can only
  depend on variables that are invariant under such rescalings. These
  are $m^2$ and $\kappa_{\ell\bar\ell}^2$ (defined in Eq.~\eqref{eq:kappa}).
  The web function admits an expansion in
  powers of the strong coupling,
  $w_{\ell \bar \ell} = \alpha_sw_{\ell \bar
    \ell}^{(1)}+\alpha_s^2w_{\ell \bar \ell}^{(2)}+\dots$, and
  \begin{equation}
    \label{eq:w1}
    w_{\ell \bar \ell}^{(1)}(m^2,m^2+\kappa^2;\varepsilon) = (4\pi)^2\frac{C_{\ell\bar\ell}}{\kappa_{\ell\bar\ell}^2}\,\mu_R^{2\epsilon}\frac{e^{\varepsilon\gamma_E}}{(4\pi)^{\varepsilon}}\,\delta(m^2)\,,
  \end{equation}
  where  $C_{\ell\bar \ell}$ is the dipole colour factor given in Eq.~\eqref{eq:Cellbell}, such that for the present case, $C_{\ell\bar \ell}=C_A/2$ if either leg is a gluon and $C_{\ell\bar \ell}=C_F-C_A/2$ otherwise.
\item the $d$-dimensional phase space, given by
  \begin{equation}
    \label{eq:dk}
    \frac{d^d k}{(2\pi)^d} = \frac{d\eta^{(\ell\bar \ell)}}{2\pi}\, dm^2 \,\frac{d\kappa^2_{\ell\bar \ell}}{(4\pi)^2}\left(\frac{\pi}{\kappa^2_{\ell\bar\ell}}\right)^\varepsilon\frac{\Gamma(1-\varepsilon)}{\Gamma(1-2\varepsilon)}\frac{d\phi}{\pi}\,(\sin^2\phi)^{-\varepsilon}\,\Theta(\sin\phi)\,. 
  \end{equation}
\end{itemize}

We extract $H_{ij}^{(1)}$ by expanding both sides of
Eq.~\eqref{eq:HQ} to order $\alpha_s$ and compare.
Evaluating the web function integral gives
\begin{multline}
  \label{eq:w-exp}
  -\alpha_s(\mu_R)  \int \frac{d^dk}{(2\pi)^d} w^{(1)}_{\ell \bar \ell}(m^2,m^2+\kappa_{\ell\bar \ell}^2;\varepsilon) \Theta\left(\frac 12\ln\frac{Q^2_{\ell\bar \ell}}{m^2+\kappa_{\ell\bar \ell}^2}-|\eta^{(\ell\bar \ell)}| \right) \Theta(Q_{\ell \bar\ell}-\kappa_{\ell\bar \ell}) \\ = -\frac{\alpha_s(\mu_R)}{2\pi}
  \left(\frac{\mu_R^2}{Q^2_{\ell\bar \ell}}\right)^\varepsilon \frac{e^{\varepsilon\gamma_E}}{\Gamma(1-\varepsilon)} \frac{C_{\ell\bar \ell}}{\varepsilon^2}\,,
\end{multline}
and the anomalous dimension integral gives
\begin{equation}
  \label{eq:gamma0-int}
  \int_0^{Q^2}\!\!\frac{dk^2}{k^2}\gamma_\ell^{(0)}\frac{\alpha_s(k^2;\varepsilon)}{2\pi} = \frac{\alpha_s(\mu_R)}{2\pi}\gamma_\ell^{(0)} \int_0^{Q^2}\!\!\frac{dk^2}{k^2}\left(\frac{\mu_R^2}{k^2}\right)^{\varepsilon} = -\frac{1}{\varepsilon}\left(\frac{\mu_R^2}{Q^2}\right)^{\varepsilon} \gamma_\ell^{(0)}\frac{\alpha_s(\mu_R)}{2\pi}\,.
\end{equation}
Using the fact that $\mathcal{V}(Q;\varepsilon)$ has the universal structure~\cite{Catani:1998bh}\footnote{Note that the general formula of~\cite{Catani:1998bh} includes also phase factors (Coulomb phases) that arise when the two legs $\ell$ and $\bar \ell$ are both incoming or outgoing. In the case of three emitters, these phase factors give rise to a finite contribution, included in $\mathcal{H}_1$.}
\begin{equation}
  \label{eq:HQ-expanded}
  \mathcal{V}(Q;\varepsilon) = 1+\frac{\alpha_s(\mu_R)}{2\pi}\left[-\frac{e^{\varepsilon\gamma_E}}{\Gamma(1-\varepsilon)}\sum_\ell\left(\frac{1}{\varepsilon^2}+\frac{1}{\varepsilon}\frac{\gamma_\ell^{(0)}}{C_\ell}\right)\sum_{\bar \ell\ne \ell} C_{\ell\bar\ell}\left(\frac{\mu_R^2}{Q_{\ell\bar \ell}^2}\right)^{\varepsilon}+\mathcal{H}_1\right]\,,
\end{equation}
and comparing the above result with the right-hand side of Eq.~(\ref{eq:HQ}),
\begin{equation}
  \label{eq:HQ-rep-expanded}
  \mathcal{V}(Q;\varepsilon) = 1+\frac{\alpha_s(\mu_R)}{2\pi}\left\{-\frac{e^{\varepsilon\gamma_E}}{\Gamma(1-\varepsilon)}\sum_\ell\left[\sum_{\bar \ell\ne \ell}\left(\frac{\mu_R^2}{Q^2_{\ell\bar\ell}}\right)^{\varepsilon}\frac{C_{\ell\bar \ell}}{\varepsilon^2}+\left(\frac{\mu_R^2}{Q^2}\right)^{\varepsilon}\frac{\gamma_\ell^{(0)}}{\varepsilon}\right]+H_1\right\}\,,
\end{equation}
we find
\begin{equation}
  \label{eq:H1-final}
  H_{ij}^{(1)}(Q,\mu_R) = \mathcal{H}_1+ \sum_\ell\frac{\gamma_\ell^{(0)}}{C_\ell}\sum_{\bar \ell\ne \ell} C_{\ell\bar \ell}\ln\left(\frac{Q^2_{\ell\bar\ell}}{Q^2}\right)\,,
\end{equation}
where the finite constant piece, $\mathcal{H}_1$, is
process dependent and determined by the one-loop hard scattering
amplitude. It must therefore be calculated separately for each hard
process.

\subsection{The NNLL Radiator}
\label{sec:radiator}

The NNLL radiator extends the NLL result derived in
Sec.~\ref{subsec:NLL-radiator} by incorporating the additional
contributions required for NNLL accuracy. Throughout this section we
retain the notation and kinematic variables introduced previously.

The radiator is obtained operationally from the singular part of the virtual corrections in
Eq.~(\ref{eq:HQ}), by cutting off hard-collinear contributions down to
the scale $Qv^{1/(a+b_\ell)}$, and including suitably defined
unresolved soft emissions. \ABquestion{I don't understand how this 
gives corrections from the veto scale and below since Eq.~\eqref{eq:Rhc-NNLL} 
clearly covers the range between the veto scale $Qv^{1/(a+b_\ell)}$ and the hard scale $Q$, rather than
below the veto scale?} The web function integral of Eq.~\eqref{eq:w-exp} and the anomalous
dimension integral of Eq.~\eqref{eq:gamma0-int}, restricted to this
region, give rise to the soft and hard-collinear pieces
respectively.

Relative to NLL accuracy, the NNLL radiator receives three additional
contributions. The first is the mass correction
$\delta R_{\ell\bar\ell}$. \ABquestion{Is this true?: This accounts for the difference
between the exact rapidity boundary and the approximation used
in the massless radiator (see below). This correction arises because unresolved collinear splittings 
are integrated inclusively.}
 The second is the second-derivative term $R_{\ell\bar\ell}''$, which originates from the azimuthal
expansion of the observable dependence in Eq.~\eqref{eq:R-ll-expan}. 
Finally, the hard-collinear contribution receives an additional contribution from the two-loop
anomalous dimension $\gamma_\ell^{(1)}$.

As at NLL, the NNLL radiator naturally decomposes into soft and
hard-collinear contributions,
\begin{equation}
  \label{eq:radiator-NNLL}
  R_{\rm NNLL}(v)=R_{\rm s}(v)+R_{\rm hc}(v),
\end{equation}
where  terms beyond NNLL accuracy are neglected.
  
  The soft contribution is given by
  \begin{equation}
  	R_{\rm s}(v)= \sum_\ell\sum_{\bar \ell\ne \ell} C_{\ell\bar \ell} \,\mathcal{R}_{\ell\bar\ell}(v)\,, 
  \end{equation}
where $\mathcal{R}_{\ell\bar\ell}(v)$ is obtained by restricting the
web-function integral to unresolved soft emissions, giving the finite
contribution
\begin{multline}
  \label{eq:Rll}
  \mathcal{R}_{\ell\bar\ell}(v) = \int\frac{d^4
    k}{(2\pi)^4}w_{\ell\bar\ell}(m^2,m^2+\kappa_{\ell\bar\ell}^2;0)
  \Theta\left(\frac 12\ln\frac{Q^2_{\ell\bar
        \ell}}{m^2+\kappa_{\ell\bar \ell}^2}-\eta^{(\ell\bar \ell)}
  \right) \Theta(Q_{\ell \bar\ell}-\kappa_{\ell\bar \ell}) \times \\
  \times \Theta(\eta^{(\ell\bar \ell)})
  \Theta\left(d^{(\ell\bar\ell)}_\ell\left(\frac{\kappa_{\ell\bar
          \ell}}{Q_{\ell\bar\ell}}\right)^{a}
    e^{-b_\ell\eta^{(\ell\bar \ell)}} g_\ell(\phi)-v\right)\,.
\end{multline}
They key difference from the NLL soft radiator of Eq.~\eqref{eq:Rll-dip}
is \ABquestion{Is the difference that we now integrate inclusively over
the multiple collinear splittings? These collinear splittings a this order are 
parameterised by the web which has a dependence on $m$?}

We note that, apart from the physical rapidity bound
rapidity, $\eta^{(\ell\bar \ell)}<\ln(Q_{\ell\bar\ell}/\sqrt{m^2+\kappa_{\ell\bar \ell}^2})$, 
  all other constraints are a matter of convention. Different choices modify the definition of the NLL function
$\FNLL$ and the NNLL correction $\delta\FNNLL$, while leaving the final resummed prediction
unchanged.

At NNLL accuracy, we can account for the $\phi$ and $m^2$ dependence of the observable 
by reorganising $\mathcal{R}_{\ell\bar\ell}(v)$ as 
\begin{equation}
\label{eq:Rll-rearranged}
  \mathcal{R}_{\ell\bar\ell}(v) = \int\frac{d\phi}{2\pi}\, R_{\ell\bar \ell}\left(\frac{v}{d_{\ell\bar \ell} \,g_\ell(\phi)},Q_{\ell\bar \ell}\right)+\delta R_{\ell\bar \ell}(v)\,,
\end{equation}
where $R_{\ell\bar\ell}(v,Q_{\ell\bar\ell})$ is the ``massless'' radiator, 
identical to the NLL radiator derived previously in
Eq.~\eqref{eq:Rll-final}. Expanding the first term to NNLL accuracy gives \ABquestion{Why is the $d_{\ell\bar \ell}g_{\ell}$ on the denominator of $v/d_{\ell\bar \ell}g_{\ell}$?}
\begin{multline}
  \label{eq:Rll-expanded}
\int\frac{d\phi}{2\pi}\, R_{\ell\bar \ell}\left(\frac{v}{d_{\ell\bar \ell} \,g_\ell(\phi)},Q_{\ell\bar \ell}\right)\simeq  
R_{\ell\bar \ell}\left(v,Q_{\ell\bar \ell}\right)\\+\langle \ln(d_{\ell\bar \ell}\,g_\ell)\rangle R'_{\ell\bar \ell}\left(v,Q_{\ell\bar \ell}\right)+\frac 12
\langle \ln^2(d_{\ell\bar \ell}\,g_\ell)\rangle R''_{\ell\bar \ell}\left(v,Q_{\ell\bar \ell}\right)\,,
\end{multline}
where the notation
$\langle\cdots\rangle$ denotes the azimuthal average defined in
Eq.~\eqref{eq:azimuthal-average} and the derivatives are defined as
$R'_{\ell\bar\ell}\equiv-vdR_{\ell\bar\ell}/dv$, and
$R''_{\ell\bar\ell}\equiv-vdR'_{\ell\bar\ell}/dv$.

The second term in Eq.~\eqref{eq:Rll-rearranged} is the ``mass correction'',
\begin{multline}
  \label{eq:dRll}
  \delta R_{\ell\bar\ell}(v) = -\int\frac{d^4
    k}{(2\pi)^4}w_{\ell\bar\ell}(m^2,m^2+\kappa_{\ell\bar\ell}^2;0)
 \Theta(Q_{\ell \bar\ell}-\kappa_{\ell\bar \ell}) \Theta(\kappa_{\ell\bar\ell}-v^{\frac1{a+b_\ell}} Q_{\ell\bar \ell}) \times \\
  \times \Theta\left(\eta^{(\ell\bar \ell)}-\frac 12\ln\frac{Q^2_{\ell\bar
        \ell}}{m^2+\kappa_{\ell\bar \ell}^2}
  \right)\Theta\left(\ln\frac{Q_{\ell\bar
        \ell}}{\kappa_{\ell\bar \ell}}-\eta^{(\ell\bar \ell)}  \right)\,,
\end{multline}
which accounts for the correct rapidity bound.

The hard-collinear radiator resums virtual corrections associated with
radiation collinear to each emitting leg. As in the NLL case (see Eq.~\eqref{eq:Rhc-NLL}),
the integration extends from the observable-dependent hard-collinear scale
$Qv^{1/(a+b_\ell)}$ up to the hard scale $Q$:
\begin{equation}
  \label{eq:Rhc-NNLL}
  R_{\rm hc}(v) =
  -\sum_\ell\int^{Q^2}_{Q^2v^{\frac{2}{a+b_\ell}}}
  \!\!\frac{dk^2}{k^2}\,\gamma_\ell\!\left(\alpha_s(k)\right)\,,
\end{equation}
where $\gamma_\ell(\alpha_s)$ is given in Eq.~\eqref{eq:gammal}. Retaining only
$\gamma_\ell^{(0)}$ recovers the NLL result of
Eq.~\eqref{eq:Rhc-NLL}, while the $\gamma_\ell^{(1)}$ term provides the new
NNLL contribution. 

Following the same procedure as at NLL, we employ the ansatz of
Eq.~\eqref{eq:ansatz} to obtain
\begin{equation}
	R_{\ell\bar\ell}(v)=-\frac{\lambda}{\alpha_s(Q)\beta_0}g_1^{(\ell)}(\lambda)-g_2^{(\ell)}(\lambda)-\mathfrak{g}_2^{(\ell)}(\lambda,Q_{\ell\bar\ell})-\frac{\alpha_s(Q)}{\pi}\left(g_3^{(\ell)}(\lambda)+\mathfrak{g}_3^{(\ell)}(\lambda,Q_{\ell\bar\ell})\right)\,,
\end{equation}
\begin{equation}
	\delta R_{\ell\bar\ell}(v)=-\frac{\alpha_s(Q)}{\pi}\delta g_3^{(\ell)}(\lambda)\,,
\end{equation}
\begin{equation}
	R_{\mathrm{hc},\ell}(v)=-h_2^{(\ell)}(\lambda)-\frac{\alpha_s(Q)}{\pi}h_3^{(\ell)}(\lambda)\,,
\end{equation}
where $g_1^{(\ell)}(\lambda)$, $g_2^{(\ell)}(\lambda)$,
$\mathfrak{g}_2^{(\ell)}(\lambda,Q_{\ell\bar\ell})$, and
$h_2^{(\ell)}(\lambda)$ are given in Eq.~\eqref{eq:RNLL-g-h}.
The additional functions, $g_3^{(\ell)}(\lambda),$ $\mathfrak{g}_3^{(\ell)}(\lambda,Q_{\ell\bar\ell})$, $\delta g_3^{(\ell)}(\lambda)$ and
$h_3^{(\ell)}(\lambda)$ are:
\begin{subequations}
\begin{align}
	&g_3^{(\ell)}(\lambda)=\left[K^{(1)}\frac{\beta_1\left(a^2(a+b_{\ell}-2\lambda)\ln\left(1-\frac{2\lambda}{a}\right)-(a+b_{\ell})^2(a-2\lambda)\ln\left(1-\frac{2\lambda}{a+b_{\ell}}\right)+6b_{\ell}\lambda^2\right)}{4\pi b_{\ell}\beta_0^3(a-2\lambda)(a+b_{\ell}-2\lambda)}\right.\notag \\[0.5ex]
	&\qquad\qquad+\frac{\left(\beta_1^2(a+b_{\ell})^2(a-2\lambda)\ln^2\left(1-\frac{2\lambda}{a+b_{\ell}}\right)-4b_{\ell}\lambda^2\left(\beta_0\beta_2+\beta_1^2\right)\right)}{4 b_{\ell}\beta_0^4(a-2\lambda)(a+b_{\ell}-2\lambda)} \notag \\[0.5ex]
	&\qquad\qquad-\frac{a\ln\left(1-\frac{2\lambda}{a}\right)\left(2\beta_0\beta_2(a-2\lambda)+a\beta_1^2\ln\left(1-\frac{2\lambda}{a}\right)+4\beta_1^2\lambda\right)}{4b_{\ell}\beta_0^4(a-2\lambda)} \notag\\[0.5ex]
	&\qquad\qquad+\frac{(a+b_{\ell})\ln\left(1-\frac{2\lambda}{a+b_{\ell}}\right)\left(\beta_0\beta_2(a+b_{\ell}-2\lambda)+2\beta_1^2\lambda\right)}{2b_{\ell}\beta_0^4(a+b_{\ell}-2\lambda)} \notag\\[0.5ex]
	&\left.\qquad\qquad -K^{(2)}\frac{2\lambda^2}{8\pi^2(a-2\lambda)(a+b_{\ell}-2\lambda)\beta_0^2} \right]\,, \\[2.5ex]
&\mathfrak{g}_3^{(\ell)}(\lambda,Q_{\ell\bar\ell})=\pi\beta_0\lambda^2\left(\frac{dg_1^{(\ell)}(\lambda)}{d\lambda}+\frac{\lambda}{2}\frac{d^2g_1^{(\ell)}(\lambda)}{d\lambda^2}\right)\ln^2\left(\frac{Q_{\ell\bar\ell}^2}{Q^2}\right)\notag \\[0.5ex]
&\qquad\qquad -\pi\beta_0\lambda\left(\frac{dg_2^{(\ell)}(\lambda)}{d\lambda}+\frac{\beta_1}{\beta_0^2}\lambda\frac{dg_1^{(\ell)}(\lambda)}{d\lambda}\right)\ln\left(\frac{Q_{\ell\bar\ell}^2}{Q^2}\right)\,, \\[2.5ex]
&\delta g_3^{(\ell)}(\lambda)=-\zeta_2\frac{\lambda}{(a+b_{\ell}-2\lambda)}\,,\\[2.5ex]
&h_3^{(\ell)}(\lambda)=-\gamma_{\ell}^{(0)}\frac{\beta_1\left((a+b_{\ell})\ln\left(1-\frac{2\lambda}{a+b_{\ell}}\right)+2\lambda\right)}{2\beta_0^2(a+b_{\ell}-2\lambda)}
+\gamma_{\ell}^{(1)}\frac{\lambda}{2\pi\beta_0(a+b_{\ell}-2\lambda)}\,,
\end{align}
\end{subequations}
where $\beta_0$, $\beta_1$, and $\beta_2$ are defined in
Eq.~\eqref{eq:beta}, while $\gamma_\ell^{(0)}$ and
$\gamma_\ell^{(1)}$ are given in
Eqs.~\eqref{eq:gammal-flav} and \eqref{eq:gammal1-flav},
respectively.

The derivatives of the NNLL radiator decompose as sums over the
emitting legs, $R'_{\mathrm{NNLL}}(\lambda)=\sum_\ell R'_{\mathrm{NNLL},\ell}(\lambda)$
and $R''_{\mathrm{NNLL}}(\lambda)=\sum_\ell R''_{\mathrm{NNLL},\ell}(\lambda)$, where
\begin{subequations}
\begin{align}
 \label{eq:RNNLL}
	&R'_{\mathrm{NNLL},\ell}(\lambda)=C_{\ell}\frac{\alpha_s(Q)}{b_{\ell}\pi^2\beta_0^2(a-2\lambda)(a+b_{\ell}-2\lambda)}
		\left[b_{\ell}\beta_0\lambda K^{(1)}-2\pi b_{\ell}\beta_1\lambda \right. \notag\\ 
		& \qquad\qquad \left.-\pi a(a+b_{\ell}-2\lambda)\beta_1\ln\left(1-\frac{2\lambda}{a}\right)+\pi(a+b_{\ell})(a-2\lambda)\beta_1\ln\left(1-\frac{2\lambda}{a+b_{\ell}}\right)\right]\,,
	 \\ \label{eq:R2NNLL}
	&R''_{\mathrm{NNLL},\ell}(\lambda)\equiv \alpha_s(Q)\beta_0\frac{dR'_{\mathrm{NLL},\ell}}{d\lambda}=2C_{\ell}\frac{\alpha_s(Q)}{\pi}\frac{1}{(a-2\lambda)(a+b_{\ell}-2\lambda)}\,.
	\end{align}
\end{subequations}
With the hard function $H^{(1)}_{ij}$ and the NNLL radiator
$R_{\mathrm{NNLL}}$ determined, we now turn to the remaining NNLL
corrections. The next section addresses the main focus of this work,
namely the derivation of initial-state radiation contributions within
the \textsc{ARES} formalism.

\section{NNLL resummation: initial-state radiation}
\label{sec:NNLL-initial}

A general treatment of initial-state radiation at NNLL is currently
missing. Results exist for one-jettiness in deep-inelastic scattering~\cite{Kang:2013ata, Kang:2013lga}
and the $N$-jettiness in hadronic collisions~\cite{Stewart:2010tn, Gaunt:2015pea}, as well as
for the transverse momentum of a colour singlet or a top-antitop pair,
and for selected observables in colour-singlet
production~\cite{vanBeekveld:2025zjh}.

We restrict ourselves to event-shape variables and, as is customary in
the \textsc{ARES} approach, we identify the relevant emissions contributing 
at NNLL: a single hard and collinear emission $k$, accompanied by an arbitrary number of soft and collinear emissions $k_1,\dots,k_n$. 
The emission $k$ carries a fraction $(1-z)/z$ of the momentum $p_\ell$ entering the hard scattering, 
\begin{equation}
  \label{eq:k||l}
  k = \frac{1-z}{z} p_\ell + k_t^{(\ell)}+ \mathcal{O}\left(\left(k_t^{(\ell)}\right)^2\right)\,.
\end{equation}
To extend \textsc{ARES} to initial-state radiation, we follow the 
same principles as for final-state radiation: 
\begin{itemize}
\item we consider one collinear emission $k$ and the corresponding virtual
  corrections, accompanied by an arbitrary
  number of soft and collinear emissions $\{k_i\}$;
\item we subtract the double counting occurring when $k$ is also soft,
since such contributions are already included at NLL in $\FNLL$.
\end{itemize}
The procedure will involve contributions that are not separately
finite in four dimensions. We use dimensional regularisation to make
each contribution finite in $4-2\varepsilon$ dimensions. For initial-state radiation, virtual corrections must include the
counterterm arising from the renormalisation of PDFs. Throughout, we assume that soft emissions coherently factorise 
from collinear radiation, i.e.\ a soft gluon sees only the total 
colour charge of harder emissions at smaller angles. 
Recently identified
coherence-violating corrections~\cite{Banfi:2025mra,Becher:2026kbr}
first enter at order $\alpha_s^4$ and require a dedicated treatment
beyond the scope of this thesis.\footnote{These corrections give an additional
contribution to $\Sigma_{ij}$, hence even with coherence effects included the structure of the NNLL cumulant stays the same.}

The parton entering the hard scattering can originate from four
possible collinear splittings. A quark entering the hard process can
arise either from an incoming quark radiating a collinear gluon,
$q\to qg$, or from an incoming gluon splitting into a
quark--antiquark pair, $g\to q\bar q$, as illustrated in
Fig.~\ref{fig:q-incoming}. Similarly, a gluon entering the hard
scattering can arise from an incoming quark emitting a
collinear gluon, $q\to gq$, or an incoming gluon splitting into two
gluons $g\to gg$, in Fig.~\ref{fig:g-incoming}.

Channels with an incoming gluon additionally give rise to spin
correlations, which generate the ISR correction
$\Delta\mathcal{F}_{\rm rec}$. This contribution vanishes for
one-jettiness because the observable has no azimuthal dependence (see Sec.~\ref{sec:hc-contributions}). 
Since it plays no role for the observable considered in this chapter, we defer the general
treatment for this correction to Appendix~\ref{app:spin-correl}.

We begin by considering the $q\to qg$ splitting, shown in
Fig.~\ref{fig:q-incoming-qqg}, and derive the corresponding NNLL
contributions explicitly. The remaining splitting channels are then
obtained by generalising this result.

\begin{figure}[H]
\centering

\begin{subfigure}{0.48\textwidth}
\centering
\begin{tikzpicture}[scale=0.85]
\begin{feynman}
  \vertex (H) at (0,0);
  \vertex (v1) at (-2.0, 1.5);
  \vertex (q0) at (-4.2, 3.0);
  \vertex (qbar) at (-4.2, -2.2);
  \vertex (gISR) at (0.6, 2.95);
  \vertex (Z)   at (3.8,  1.6);
  \vertex (gout) at (3.8, -1.6);
  \diagram*{
    (q0)   -- [fermion] (v1),
    (v1)   -- [fermion, edge label={$z$}] (H),
    (v1)   -- [gluon]                     (gISR),
    (qbar) -- [anti fermion]              (H),
    (H)    -- [boson]                     (Z),
    (H)    -- [gluon]                     (gout),
  };
  \draw[fill=white] (H) circle (0.8cm) node {$\mathcal{H}_{q\bar{q}}$};
\end{feynman}
\node[above left]  at (-4.2, 3.0)  {$q$};
\node[above right] at (-2.1, 2.3)  {$1-z$};
\node[below left]  at (-4.2, -2.2) {$\bar{q}$};
\node[above right] at (3.8,  1.6)  {$Z$};
\node[below right] at (3.8, -1.6)  {$g$};
\end{tikzpicture}
\caption{$q \to qg$ splitting}
\label{fig:q-incoming-qqg}
\end{subfigure}
\hfill
\begin{subfigure}{0.48\textwidth}
\centering
\begin{tikzpicture}[scale=0.85]
\begin{feynman}
  \vertex (H) at (0,0);
  \vertex (v1) at (-2.0, 1.5);
  \vertex (q0) at (-4.2, 3.0);
  \vertex (qbar) at (-4.2, -2.2);
  \vertex (gISR) at (0.6, 2.95);
  \vertex (Z)   at (3.8,  1.6);
  \vertex (gout) at (3.8, -1.6);
  \diagram*{
    (q0)   -- [gluon] (v1),
    (v1)   -- [fermion, edge label={$z$}] (H),
    (v1)   -- [anti fermion]              (gISR),
    (qbar) -- [anti fermion]              (H),
    (H)    -- [boson]                     (Z),
    (H)    -- [gluon]                     (gout),
  };
  \draw[fill=white] (H) circle (0.8cm) node {$\mathcal{H}_{q\bar{q}}$};
\end{feynman}
\node[above left]  at (-4.2, 3.0)  {$g$};
\node[above right] at (-2.1, 2.3)  {$1-z$};
\node[below left]  at (-4.2, -2.2) {$\bar{q}$};
\node[above right] at (3.8,  1.6)  {$Z$};
\node[below right] at (3.8, -1.6)  {$g$};
\end{tikzpicture}
\caption{$g \to q\bar{q}$ splitting}
\label{fig:q-incoming-gqq}
\end{subfigure}

\caption{Initial-state radiation contributions to the $q\bar{q}$ channel
of $Z$+jet production, from the $q \to qg$ splitting (left) and the $g \to q\bar{q}$ splitting (right).  
The hard scattering is denoted by $\mathcal{H}_{q\bar{q}}$,
where the subscripts specify the incoming partons. The momentum
fractions $z$ and $1-z$ correspond to the daughter and emitted partons,
respectively, relative to the parent parton carrying momentum
$x_\ell/z$.}

\label{fig:q-incoming}
\end{figure}

\begin{figure}[H]
\centering

\begin{subfigure}{0.48\textwidth}
\centering
\begin{tikzpicture}[scale=0.85]
\begin{feynman}
  \vertex (H) at (0,0);
  \vertex (v1) at (-2.0, 1.5);
  \vertex (q0) at (-4.2, 3.0);
  \vertex (qbar) at (-4.2, -2.2);
  \vertex (gISR) at (0.6, 2.95);
  \vertex (Z)   at (3.8,  1.6);
  \vertex (gout) at (3.8, -1.6);
  \diagram*{
    (q0)   -- [fermion] (v1),
    (v1)   -- [gluon,   edge label={$z$}] (H),
    (v1)   -- [fermion]                   (gISR),
    (qbar) -- [fermion]              (H),
    (H)    -- [boson]                     (Z),
    (H)    -- [fermion]                     (gout),
  };
  \draw[fill=white] (H) circle (0.8cm) node {$\mathcal{H}_{gq}$};
\end{feynman}
\node[above left]  at (-4.2, 3.0)  {$q$};
\node[above right] at (-2.1, 2.3)  {$1-z$};
\node[below left]  at (-4.2, -2.2) {$q$};
\node[above right] at (3.8,  1.6)  {$Z$};
\node[below right] at (3.8, -1.6)  {$q$};
\end{tikzpicture}
\caption{$q \to gq$ splitting}
\label{fig:g-incoming-qgq}
\end{subfigure}
\hfill
\begin{subfigure}{0.48\textwidth}
\centering
\begin{tikzpicture}[scale=0.85]
\begin{feynman}
  \vertex (H) at (0,0);
  \vertex (v1) at (-2.0, 1.5);
  \vertex (q0) at (-4.2, 3.0);
  \vertex (qbar) at (-4.2, -2.2);
  \vertex (gISR) at (0.6, 2.95);
  \vertex (Z)   at (3.8,  1.6);
  \vertex (gout) at (3.8, -1.6);
  \diagram*{
    (q0)   -- [gluon] (v1),
    (v1)   -- [gluon, edge label={$z$}] (H),
    (v1)   -- [gluon]                   (gISR),
    (qbar) -- [fermion]            (H),
    (H)    -- [boson]                   (Z),
    (H)    -- [fermion]                   (gout),
  };
  \draw[fill=white] (H) circle (0.8cm) node {$\mathcal{H}_{gq}$};
\end{feynman}
\node[above left]  at (-4.2, 3.0)  {$g$};
\node[above right] at (-2.1, 2.3)  {$1-z$};
\node[below left]  at (-4.2, -2.2) {$q$};
\node[above right] at (3.8,  1.6)  {$Z$};
\node[below right] at (3.8, -1.6)  {$q$};
\end{tikzpicture}
\caption{$g \to gg$ splitting}
\label{fig:g-incoming-ggg}
\end{subfigure}

\caption{Initial-state radiation contributions to the $gq$ channel of
$Z$+jet production, from the $q \to gq$ splitting (left) and the $g \to
gg$ splitting (right). The hard scattering is denoted by
$\mathcal{H}_{gq}$, where the subscripts indicate the incoming
partons. The momentum fractions $z$ and $1-z$ correspond to the
daughter and emitted partons, respectively, relative to the parent
parton carrying momentum $x_\ell/z$.}
\label{fig:g-incoming}
\end{figure}

We consider an incoming quark carrying a fraction $x_{\ell}/z$ of the proton momentum, which emits a collinear gluon carrying a fraction $1-z$ of its momentum. After the emission, the quark enters the hard scattering
with a fraction $x$ of the proton momentum. The contribution of this emission to $\Sigma_{ij}(\{\tilde p\},v)$ defines a function $\mathcal{F}_{qq}$, which is reorganised to give various NNLL contributions. Our starting point is
\begin{equation}
  \label{eq:FISR-start}
  \begin{aligned}
    &\mathcal{F}_{qq} = \int \dZ \int_0^{2\pi}\frac{d\phi}{2\pi}\left\{\int_0^{Q^2}\frac{dk_t^2}{k_t^2}\left(\frac{\mu_R^2}{k_t^2}\right)^{\varepsilon} \frac{\alpha_s(k_t)}{2\pi}\right. \\
    & \quad \times \left[\int_{x_\ell}^1 \frac{dz}{z}\hat P_{qq}(z;\varepsilon)\Theta\left(1-\frac{V_{\mathrm{hc}}(\{\tilde{p}\},k,\{k_i\})}{v}\right) \frac{f_q(x_\ell/z,\mu_F)}{f_q(x_\ell,\mu_F)}- \int_0^1 dz \frac{2C_\ell}{1-z}\Theta\left(1-\frac{V_{\mathrm{sc}}(\{\tilde{p}\},k,\{k_i\}}{v}\right)\right] \\
   & \quad +\left[-\gamma_\ell^{(0)}\int_0^{Q^2v^{2/(a+b_\ell)}}\frac{dk^2}{k^2}\left(\frac{\mu_R^2}{k^2}\right)^{\varepsilon}\frac{\alpha_s(k)}{2\pi}
     +\frac{1}{\varepsilon} \left(\frac{\mu_R^2}{\mu_F^2}\right)^{\varepsilon} \frac{\alpha_s(\mu_R)}{2\pi} \int_{x_\ell}^1 \frac{dz}{z} P_{qq}(z) \frac{f_q(x_\ell/z,\mu_F)}{f_q(x_\ell,\mu_F)}\right]  \\ 
     & \quad \times \left. \Theta\left(1-\frac{V_{\mathrm{sc}}(\{\tilde{p}\},\{k_i\}}{v}\right)\right\} \,.
\end{aligned}
\end{equation}
where we have introduced the regularised Altarelli-Parisi (AP) splitting function
\begin{equation}
  \label{eq:regularised-Pij}
  \hat P_{ik}(z;\varepsilon) = \hat P_{ik}(z;0)+\varepsilon P^{(\varepsilon)}_{ik}(z)\,, \qquad k=q,g\,. 
\end{equation}
In the present case,
\begin{equation}
  \label{eq:Pqq}
  \hat P_{qq}(z;0)  = C_F \frac{1+z^2}{1-z}\,, \qquad \hat P_{qq}^{(\varepsilon)}(z)=-C_F (1-z)\,,
\end{equation}
and $P_{qq}(z) = (\hat{P}_{qq}(z;0))_+$ is the full AP splitting function.

We now comment on the physical meaning of the ingredients within Eq.~\eqref{eq:FISR-start}.
The second line represents the splitting of a
quark emitting a collinear gluon $k$ in $4-2\varepsilon$ dimensions,
minus the contribution when the gluon is soft, already accounted for
in $\mathcal{F}_{\rm NLL}$. The third line contains virtual
corrections, as well as the collinear counterterm needed to renormalise
the parton distribution function. The parameter $b_\ell$ corresponds
to the leg identified by the quark entering the hard process.  The
expression in Eq.~\eqref{eq:FISR-start} contains single logarithms of
the form $\alpha_s\ln[Qv^{1/(a+b_\ell)}/\mu_R]$ and
$\alpha_s\ln[Qv^{1/(a+b_\ell)}/\mu_F]$. These logarithms can be
resummed in the PDF ratio in Eq.~\eqref{eq:Sigma-NNLL}
by fixing $\mu_R=\mu_F=Qv^{1/(a+b_\ell)}$.
At NNLL, we can set
$\alpha_s(k)=\alpha_s(k_t)=\alpha_s(Qv^{1/(a+b_\ell)})$ and obtain a
finite expression, up to subleading terms of order $\alpha_s^2$.

In order to compute $\mathcal{F}_{qq}$ in four dimensions, we add and subtract
a counterterm, $\mathrm{CT}$, which has the following form:
\begin{multline}
    \mathrm{CT}=\int \dZ \frac{\alpha_s(Qv^{1/(a+b_\ell)})}{2\pi} \int\frac{dk_t^2}{k_t^2}\left(\frac{\mu_R^2}{k_t^2}\right)^{\varepsilon} \int_0^{2\pi}\frac{d\phi}{2\pi} \times \\
    \times \left[\int_{x_\ell}^1\frac{dz}{z}\hat P_{qq}(z;\varepsilon)\frac{f_q(x_\ell/z,Qv^{\frac{1}{a+b_\ell}})}{f_q(x_\ell,Qv^{\frac{1}{a+b_\ell}})}-\int_0^1dz\frac{2 C_\ell}{1-z}\right]\times \\ \times \Theta\left(1-\frac{V_{\mathrm{sc}}(\{\tilde{p}\},k)}{v}\right)\Theta\left(1-\frac{V_{\mathrm{sc}}(\{\tilde{p}\},\{k_i\}}{v}\right)\,.
  \end{multline}
  Adding and subtracting $\mathrm{CT}$ recasts $\mathcal{F}_{qq}$ in
the form of Eq.~\eqref{eq:Sigma-NNLL}:
\begin{equation}
  \label{eq:FISR-split}
  \mathcal{F}_{qq} = \frac{\alpha_s(Q)}{\pi}\left( \delta \mathcal{F}_{\rm rec,qq}+\delta\mathcal{F}_{\rm hc,qq}\right)+ \frac{\alpha_s(Q v^{1/(a+b_\ell)})}{2\pi} C^{(1)}_{\rm hc,qq} \mathcal{F}_{\rm NLL}\,.
\end{equation}
The coefficient $C^{(1)}_{\rm hc,qq}$ plays the same role as an integrated counterterm in fixed-order calculations, 
\ABquestion{How exactly does  $C^{(1)}_{\rm hc,qq}$  play the role of an integrated counterterm?} whilst the sum $\delta \mathcal{F}_{\rm rec,qq}+\delta\mathcal{F}_{\rm hc,qq}$ is the real emission contribution after $\mathrm{CT}$ has been
subtracted, leaving a finite integrand in four dimensions. We now proceed to the detailed derivation of Eq.~\eqref{eq:FISR-split}. 

The coefficient $C^{(1)}_{\mathrm{hc, qq}}$ is obtained by adding $\mathrm{CT}$ to the last two lines of Eq.~\eqref{eq:FISR-start}:
\begin{equation}
\begin{aligned}
  & C^{(1)}_{\rm hc,qq} \mathcal{F}_{\rm NLL}(\lambda)= \int \dZ \int_0^{2\pi}\frac{d\phi}{2\pi}\,\Theta\left(1-\frac{V_{\mathrm{sc}}(\{\tilde{p}\},\{k_i\})}{v}\right) \times \\ & \times \left\{ -\gamma_\ell^{(0)}\int_0^{Q^2v^{2/(a+b_\ell)}}\frac{dk^2}{k^2}\left(\frac{Q^2v^{2/(a+b_\ell)}}{k^2}\right)^{\varepsilon}+\frac{1}{\varepsilon} \int_{x_\ell}^1 \frac{dz}{z} P_{qq}(z) \frac{f_q(x_\ell/z,Q v^{1/(a+b_\ell)})}{f_q(x_\ell,Q v^{1/(a+b_\ell)})} \right. \\ 
   &\left. +\int_0^{Q^2}\frac{dk_t^2}{k_t^2}\left(\frac{Q^2v^{2/(a+b_\ell)}}{k_t^2}\right)^{\varepsilon} \left[\int_{x_\ell}^1\frac{dz}{z} \hat P_{qq}(z;\varepsilon)\frac{f_q(x_\ell/z,Qv^{\frac{1}{a+b_\ell}})}{f_q(x_\ell,Qv^{\frac{1}{a+b_\ell}})}-\int_0^1dz\frac{2 C_\ell}{1-z}\right]\Theta\left(1-\frac{V_{\mathrm{sc}}(\{\tilde{p}\},k)}{v}\right)\right\}\,.
\end{aligned}
\end{equation}
By comparing the first term on the right-hand-side of the above
equation to Eq.~\eqref{eq:FNLL-formula}, we recognise $\FNLL(\lambda)$. Using the parametrisation of
$V_{\mathrm{sc}}(\{\tilde{p}\},k)$ in terms of the splitting variables
\begin{equation}
\label{eq:Vsc-param}
V_{\mathrm{sc}}(\{\tilde{p}\},k)=d_\ell g_\ell(\phi) \left(\frac{Q}{2 E_\ell}\right)^{b_\ell}\left(\frac{k_t}{Q}\right)^{a+b_\ell}\frac{1}{(1-z)^{b_\ell}}\,,
\end{equation}
together with $\gamma_\ell^{(0)}=-3/2 C_F$, we find 
\begin{multline}
C^{(1)}_{\rm hc, qq}=  \int_{x_\ell}^1\frac{dz}{z}\,\frac{f_q(x_\ell/z,Qv^{\frac{1}{a+b_\ell}})}{f_q(x_\ell,Qv^{\frac{1}{a+b_\ell}})}\times \\ \times
C_F\left[\frac{2 }{a+b_\ell}\left(b_\ell\left(\frac{\ln(1-z)}{1-z}\right)_+-\frac{1}{(1-z)_+}
        \left(\langle\ln d_\ell g_\ell\rangle-b_\ell \ln\frac{2E_\ell}{Q}\right)\right)(1+z^2)+(1-z)\right]\,.
  \end{multline}
  Using the regularised AP function in Eq.~\eqref{eq:Pqq}, the above expression can be written in the more compact form
  \begin{multline}
    \label{eq:C1qq-final}
C^{(1)}_{\rm hc, qq}=  \int_{x_\ell}^1\frac{dz}{z}\,\frac{f_q(x_\ell/z,Qv^{\frac{1}{a+b_\ell}})}{f_q(x_\ell,Qv^{\frac{1}{a+b_\ell}})}\times \\ \times
\left[\frac{2 }{a+b_\ell}\left(b_\ell\left(\frac{\ln(1-z)}{1-z}\right)_+-\frac{1}{(1-z)_+}
        \left(\langle\ln d_\ell g_\ell\rangle-b_\ell \ln\frac{2E_\ell}{Q}\right)\right)(1-z)\hat P_{qq}(z;0)-P_{qq}^{(\varepsilon)}(z)\right]\,.
  \end{multline}
  Referring back to Eq.~\eqref{eq:FISR-split}, we now evaluate the remaining contribution 
$\frac{\alpha_s}{\pi}\left(\delta \mathcal{F}_{\mathrm{rec,qq}}+\delta \mathcal{F}_{\mathrm{hc,qq}}\right)$. 
In analogy with the extraction of $C^{(1)}_{\mathrm{hc},qq}$, this term is obtained by subtracting the counterterm $\mathrm{CT}$ from the first two lines of Eq.~\eqref{eq:FISR-start}, yielding
\begin{equation}
\label{eq:dFISR+CT}
\begin{aligned}
  \frac{\alpha_s(Q)}{\pi}&\left(\delta \mathcal{F}_{\mathrm{rec,qq}}+\delta \mathcal{F}_{\mathrm{hc,qq}}\right)=\int \dZ
  \frac{\alpha_s(Qv^{\frac{1}{a+b_\ell}})}{2\pi} \int_0^{Q^2}\frac{dk_t^2}{k_t^2}
\left(\frac{Q^2v^{\frac{2}{a+b_\ell}}}{k_t^2}\right)^{\varepsilon} \int_0^{2\pi}\frac{d\phi}{2\pi}
\times \\
&\times \left\{\left[\int_{x_\ell}^1 \frac{dz}{z}
\hat P_{qq}(z;\epsilon)
\Theta\left(1-\frac{V_{\mathrm{hc}}(\{\tilde{p}\},k,\{k_i\})}{v}\right)
\frac{f_q(x_\ell/z,Qv^{\frac{1}{a+b_\ell}})}{f_q(x_\ell,Qv^{\frac{1}{a+b_\ell}})}+\right.\right. \\
&\left.
  - \int_0^1 dz \frac{2C_\ell}{1-z}
\Theta\left(1-\frac{V_{\mathrm{sc}}(\{\tilde{p}\},k,\{k_i\})}{v}\right)\right] + \\
&-\left[\int_{x_\ell}^1\frac{dz}{z} \hat P_{qq}(z;\epsilon)
\frac{f_q(x_\ell/z,Qv^{\frac{1}{a+b_\ell}})}{f_q(x_\ell,Qv^{\frac{1}{a+b_\ell}})}
-\int_0^1dz\frac{2 C_\ell}{1-z}\right] \times\\
&\times \left.\Theta\left(1-\frac{V_{\mathrm{sc}}(\{\tilde{p}\},k)}{v}\right)
\Theta\left(1-\frac{V_{\mathrm{sc}}(\{\tilde{p}\},\{k_i\})}{v}\right)
\right\}\,.
\end{aligned}
\end{equation}
To simplify the above expression, we introduce the rescaled parameter $
\zeta \equiv V_{\mathrm{sc}}(\{\tilde{p}\},k)/v$, and change variables using the parametrisation of
$V_{\mathrm{sc}}(\{\tilde{p}\},k)$ in Eq.~\eqref{eq:Vsc-param}. This gives
\begin{equation}
	\frac{dk_t^2}{k_t^2}\simeq\frac{2}{\left(a+b_\ell\right)}\frac{d\zeta}{\zeta}\,.
\end{equation}
Using this substitution and taking the limit $\varepsilon \to 0$, Eq.~\eqref{eq:dFISR+CT} simplifies to
\begin{equation}
  \begin{aligned}
    \label{eq:dFrec+dFhq}
\frac{\alpha_s(Q)}{\pi}&\left(\delta \mathcal{F}_{\mathrm{rec,qq}}+\delta \mathcal{F}_{\mathrm{hc,qq}}\right)=\int \dZ \frac{\alpha_s\left(Qv^{1/(a+b_\ell)}\right)}{\pi\left(a+b_\ell\right)}\left\{\int_0^\infty\frac{d\zeta}{\zeta}\int_0^{2\pi}\frac{d\phi}{2\pi}\times\right. \\
&\times \left(\int_{x_\ell}^1 \frac{dz}{z} \hat P_{qq}(z;0)\,\frac{f_q(x_\ell/z,Qv^{1/(a+b_\ell)})}{f_q(x_\ell,Qv^{1/(a+b_\ell)})}
\left[\Theta\left(1-\frac{V_{\mathrm{hc}}(\{\tilde{p}\},k,\{k_i\})}{v}\right)+\right.\right.\\
&\left.-\Theta(1-\zeta)\Theta\left(1-\frac{V_{\mathrm{sc}}(\{\tilde{p}\},\{k_i\})}{v}\right)\right] +\\
&\left.\left.-\int_0^1 dz \frac{2C_\ell}{1-z}\left[\Theta\left(1-\frac{V_{\mathrm{sc}}(\{\tilde{p}\},k,\{k_i\})}{v}\right) -\Theta(1-\zeta)\Theta\left(1-\frac{V_{\mathrm{sc}}(\{\tilde{p}\},\{k_i\})}{v}\right)\right]\right)\right\}\,.
\end{aligned}
\end{equation}
We observe that the second and third line give two contributions that
are separately finite in the collinear limit $\zeta\to 0$. However, in
order to take the limit
$z\to 1$ and obtain a finite results, one needs to compute the expression in
Eq.~\eqref{eq:dFrec+dFhq} as a whole. A further simplification can be achieved by adding and subtracting the following term
\begin{multline}
  \label{eq:dFhc-sub}
  \int \dZ \frac{\alpha_s\left(Q v^{1/(a+b_\ell)}\right)}{\pi\left(a+b_\ell\right)}\int_0^\infty\frac{d\zeta}{\zeta}\int_0^{2\pi}\frac{d\phi}{2\pi}
 \int_{x_\ell}^1 \frac{dz}{z}  \hat P_{qq}(z;0)\, \frac{f_q(x_\ell/z,Qv^{1/(a+b_\ell)})}{f_q(x_\ell,Qv^{1/(a+b_\ell)})}\times\\
 \times\Theta\left(1-\frac{V_{\mathrm{sc}}(\{\tilde{p}\},k,\{k_i\})}{v}\right)\,.
\end{multline}
This gives two contributions, namely $\delta \mathcal{F}_{\mathrm{rec,qq}}$ and $\delta \mathcal{F}_{\mathrm{hc,qq}}$:
\begin{equation}
\label{eq:dFrec-qq}
\begin{aligned}
\delta \mathcal{F}_{\mathrm{rec,qq}} = \frac{\alpha_s\left(Q v^{1/(a+b_\ell)}\right)}{\alpha_s(Q)\left(a+b_\ell\right)}\int_0^\infty\frac{d\zeta}{\zeta}&\int_0^{2\pi}\frac{d\phi}{2\pi} \int \dZ\int_{x_\ell}^1 \frac{dz}{z} \hat P_{qq}(z;0)\,\frac{f_q(x_\ell/z,Qv^{1/(a+b_\ell)})}{f_q(x_\ell,Qv^{1/(a+b_\ell)})} \\
\times  
&\left[\Theta\left(1-\frac{V_{\mathrm{hc}}(\{\tilde{p}\},k,\{k_i\})}{v}\right)-\Theta\left(1-\frac{V_{\mathrm{sc}}(\{\tilde{p}\},k,\{k_i\})}{v}\right)\right]\,,
\end{aligned}
\end{equation}
\begin{equation}
\label{eq:dFhc-qq}
\begin{aligned}
\delta \mathcal{F}_{\mathrm{hc,qq}} &= \frac{\alpha_s\left(v^{1/(a+b_\ell)}\right)}{\alpha_s(Q)\left(a+b_\ell\right)}\left(\int_{x_\ell}^1 \frac{dz}{z}\frac{f_q(x_\ell/z,Qv^{1/(a+b_\ell)})}{f_q(x_\ell,Qv^{1/(a+b_\ell)})} \left(P_{qq}(z)+\gamma^{(0)}_\ell\delta(1-z)\right)\right)\int_0^\infty\frac{d\zeta}{\zeta}\int_0^{2\pi}\frac{d\phi}{2\pi}\times \\
&\times \int \dZ   \left[\Theta\left(1-\frac{V_{\mathrm{sc}}(\{\tilde{p}\},k,\{k_i\})}{v}\right) -\Theta(1-\zeta)\Theta\left(1-\frac{V_{\mathrm{sc}}(\{\tilde{p}\},\{k_i\})}{v}\right)\right]\,.
\end{aligned}
\end{equation} 
 This channel $q\to qg$ is an instance of a
general structure. The coefficient $C^{(1)}_{\mathrm{hc}, qq}$ can be written in the general form,
  \begin{equation}
    \label{eq:Cell-ISR}
    C_{\mathrm{hc},\ell}^{(1)} = \sum_k  \int_{x_\ell}^1\frac{dz}{z}\,C_{ik}^{(1)}(z) \frac{f_k(x_\ell/z,Qv^{\frac{1}{a+b_\ell}})}{f_i(x_\ell,Qv^{\frac{1}{a+b_\ell}})}\,,
  \end{equation}
  where $i$ is the flavour of leg $\ell$. 
  \ABquestion{I think there is a notation issue here -- I believe we 
  need to say $i$ is the flavour of $\ell=1$ or instead make the subscript 
  $i_{\ell}$ not $i$ as the channel is $ij$.}
  The coefficients $C_{ik}^{(1)}(z)$ are given in terms of the splitting functions in $4-2\varepsilon$ dimensions
  \begin{equation}
    \label{eq:Pij-general}
    \hat P_{ik}(z;\varepsilon)= \hat P_{ik}(z;0)+\varepsilon P_{ik}^{(\varepsilon)}(z)\,,
  \end{equation}
  and have the following expression
  \begin{multline}
    \label{eq:Cij-ISR}
    C_{ik}^{(1)}(z) = -\hat P_{ik}^{(\varepsilon)}(z) \\ +\frac{2 }{a+b_\ell}\left(b_\ell\left(\frac{\ln(1-z)}{1-z}\right)_+-\frac{1}{(1-z)_+}
      \left(\langle\ln d_\ell g_\ell\rangle-b_\ell \ln\frac{2E_\ell}{Q}\right)\right)(1-z)\hat P_{ik}(z;0)\,.
  \end{multline}
  The explicit expressions of the splitting
  functions are
  \begin{subequations}
  	\begin{align}
	& \hat P_{qq}(z;\varepsilon)=C_F \frac{1+z^2}{1-z}-\varepsilon\,C_F (1-z)\,,\\
	& \hat P_{qg}(z;\varepsilon)=T_R \left[z^2+(1-z)^2\right]-\varepsilon\,T_R\left[2z(1-z)\right]\,, \\
	  \label{eq:Pgq-epsilon}
	& \hat P_{gq}(z;\varepsilon) =C_F \frac{1 + (1-z)^2}{z} - \varepsilon\,C_Fz \,, \\
	 \label{eq:Pgg-epsilon}
	& \hat P_{gg}(z;\varepsilon) =2C_A\left[\frac{z}{1-z} + \frac{1-z}{z} + z(1-z)\right]\,.
	\end{align}
  \end{subequations}
  For the final-state leg ($\ell=3$), $C^{(1)}_{{\rm hc},3}$ depends
  on the flavour of the outgoing parton and can be obtained by
  the replacement
  \begin{equation}
    \label{eq:C-recovery}
    \frac{1}{z}\frac{f_k(x_\ell/z,Qv^{\frac{1}{a+b_\ell}})}{f_i(x_\ell,Qv^{\frac{1}{a+b_\ell}})}\Theta(z-x_\ell)\to 1
  \end{equation}
  and integrating over $z$. This gives exactly
  the same results as in Ref.~\cite{Arpino:2019ozn},
  such that for a quark or antiquark,
  \begin{equation}
    \label{eq:C1hc-q}
    C_{\mathrm{hc},3}^{(1)}=C_F\left(\frac{7}{2}\frac{b_3}{a+b_3}+\frac{3}{a+b_3}\left(\langle \ln d_3 g_3(\phi)\rangle-b_3\ln\frac{2E_3}{Q}\right)+\frac 12\right)\,,
  \end{equation}
  while for a gluon
  \begin{equation}
    \label{eq:C1hc-g}
    \begin{split}
    C_{\mathrm{hc},3}^{(1)}& =\left(\frac{67}{18}C_A-\frac{13}{9}T_R n_f\right)\frac{b_3}{a+b_3}\\ & +\frac{1}{a+b_3}\left(\frac{11}{3}C_A-\frac{4}{3}T_R n_f\right)\left(\langle \ln d_3 g_3(\phi)\rangle-b_3\ln\frac{2E_3}{Q}\right)+\frac{1}{3}T_R n_f\,.
    \end{split}
  \end{equation}

Similarly, we find that $\delta \mathcal{F}_{\rm hc}$ can be written generally as
 \begin{multline}
   \label{eq:dFhc-ISR}
   \delta\mathcal{F}_{\rm hc}(\lambda)=\sum_\ell \frac{\alpha_s\left(Qv^{1/(a+b_\ell)}\right)}{\alpha_s(Q)\left(a+b_\ell\right)}B_\ell \int_0^{\infty}\frac{d\zeta}{\zeta}\int_0^{2\pi}\frac{d\phi}{2\pi}
 \int \dZ  \times\\ \times  \left[\Theta\left(1-\frac{V_{\mathrm{sc}}(\{\tilde{p}\},k,\{k_i\})}{v}\right) -\Theta(1-\zeta)\Theta\left(1-\frac{V_{\mathrm{sc}}(\{\tilde{p}\},\{k_i\})}{v}\right)\right]\,.
 \end{multline}
 For initial-state radiation, $B_\ell$ is given by
 \begin{equation}
   \label{eq:Bell-ISR}
B_\ell=\sum_k\int_{x_{\ell}}^1 \frac{dz}{z} \left(P_{ik}(z)-\gamma_\ell^{(0)}\delta_{ik}\right) \frac{f_k(x_{\ell}/z,Qv^{1/(a+b_\ell)})}{f_i(x_\ell,Qv^{1/(a+b_\ell)})}  \,,
 \end{equation}
 where $P_{ik}(z)$ are the full Altarelli-Parisi splitting functions,
 \begin{subequations}
 \label{eq:full-AP}
\begin{align}
	&P_{qq}(z)=C_F\left[\frac{1+z^2}{1-z}\right]_+\,,\\
	&P_{qg}(z)=T_R\left[z^2+(1-z)^2\right]\,,\\
	&P_{gq}(z)=C_F\left[\frac{1+(1-z)^2}{z}\right]\,,\\
	&P_{gg}(z)=2C_A\left[\frac{z}{(1-z)_+}+\frac{1-z}{z}+z(1-z)\right]+2\pi\beta_0\delta(1-z)\,.
\end{align}
 \end{subequations}
  For final-state radiation, the same substitution can be
  used as in Eq.~\eqref{eq:C-recovery}, such that $B_\ell=\gamma_\ell^{(0)}$,
 where 
 \begin{equation}
  \label{eq:gammal-flav}
  \gamma_\ell^{(0)} = \left\{
    \begin{split}
      & \frac{3}{2}C_F\,,\qquad\text{if $\ell$ is a quark,} \\
      & 2\pi \beta_0 \,,\qquad\text{if $\ell$ is a gluon}\,.
    \end{split}
    \right.
\end{equation}

The recoil correction, $\delta\mathcal{F}_{\rm rec}$ can be written in the general form,
 \begin{multline}
   \label{eq:dFrec-ISR}
   \delta\mathcal{F}_{\rm rec}(\lambda)=\sum_\ell \frac{\alpha_s\left(Qv^{1/(a+b_\ell)}\right)}{\alpha_s(Q)\left(a+b_\ell\right)}\int_0^1 dz P_\ell(z) \int_0^{\infty}\frac{d\zeta}{\zeta}\int_0^{2\pi}\frac{d\phi}{2\pi}
 \int \dZ  \times\\ \times  \left[\Theta\left(1-\frac{V_{\mathrm{hc}}(\{\tilde{p}\},k,\{k_i\})}{v}\right) -\Theta\left(1-\frac{V_{\mathrm{sc}}(\{\tilde{p}\},k,\{k_i\})}{v}\right)\right]\,.
\end{multline}
For an initial-state leg,
\begin{equation}
  \label{eq:Pell-ISR}
  P_\ell(z)=\sum_k \hat P_{ik}(z;0) \frac{f_k(x_\ell/z,Qv^{1/(a+b_\ell)})}{z f_i(x_\ell,Qv^{1/(a+b_\ell)})} \Theta(z-x_\ell)\,,
\end{equation}
where $ \hat P_{ik}(z;0)$ are defined in Eq.~\eqref{eq:Pij-general}.
For a final-state leg, 
\begin{equation}
  \label{eq:Pell-FSR}
  P_\ell(z)=\left\{
    \begin{split}
      &C_F \frac{1+z^2}{1-z}\,,\qquad\qquad\qquad\qquad\qquad\qquad\qquad\qquad\quad\text{if $\ell$ is a quark,}\\
      &C_A\left(\frac{2z}{1-z}+z(1-z)\right)+T_R\left[z^2+(1-z)^2\right]\,,\qquad\text{if $\ell$ is a gluon}\,.
    \end{split}
  \right.
\end{equation}

\section{ NNLL resummation: one-jettiness}
\label{sec:one-jettiness}
We now apply the general NNLL formalism to one-jettiness $\tau_1$,
introduced in Sec.~\ref{sec:setup}. One-jettiness is an rIRC-safe,
continuously global observable that vanishes at Born level, and satisfies the scaling conditions for an arbitrary number of soft
and/or collinear emissions discussed in Sec.~\ref{sec:obs-req}. It
therefore belongs to the class of observables to which the
\textsc{ARES} formalism applies.

We begin by considering the behaviour of the observable in the presence of a
single soft emission collinear to leg $\ell$. By definition,
 $k\cdot p_\ell/E_\ell=k_t^{(\ell)} e^{-\eta^{(\ell)}}$, where
  $k_t^{(\ell)}$ is the relative transverse momentum with respect to
  the emitting leg $\ell$ and $\eta^{(\ell)}$ its rapidity with respect to the same leg. 
  Using $\tau_1 = \mathcal{T}_1/M$ with $\mathcal{T}_1$ defined in
Eq.~\eqref{eq:T1}, the soft-collinear limit takes the same form for each leg:
    \begin{equation}
    \label{eq:tau1-sc-1emsn}
    \tau_1(\{\tilde p\},k)\simeq \tau^{\rm sc}_1(\{\tilde p\},k) =  \frac{2E_\ell}{Q_\ell}\frac{k_t^{(\ell)}}{M} e^{-\eta^{(\ell)}}\,.
  \end{equation}
  From Eq.~\eqref{eq:tau1-sc-1emsn}  we immediately identify the single-emission
parameters,
\begin{equation}
  \label{eq:tau1-1em-props}
  a=1\,, \quad b_\ell=1\,, \quad d_\ell=\frac{Q}{M}\frac{2E_{\ell}}{Q_{\ell}} \,, \quad g_\ell(\phi) = 1\,, \qquad \ell=1,2,3\,.
\end{equation}
Substituting these into the general expressions, 
Eqs.~\eqref{eq:H1-final},~\eqref{eq:radiator-NNLL},
and~\eqref{eq:Cell-ISR}, determines the hard coefficient
$H_{ij}^{(1)}$, the NNLL radiator $R_{\rm NNLL}$, and the
hard-collinear coefficients $C_{\rm hc,\ell}^{(1)}$,
respectively. Since one-jettiness is additive, the NLL function
$\mathcal{F}_{\rm NLL}$ is given directly by the additive expression in
Eq.~\eqref{eq:FNLL-add-final}. The remaining NNLL corrections are those 
entering Eq.~\eqref{eq:dFNNLL}. For additive observables these expressions
simplify considerably, and the corresponding master formulae are collected in
Appendix~\ref{ap:add-res-NNLL}. In the
following subsections we evaluate each contribution separately,
organising the discussion according to the kinematic region from which
the correction originates.

\subsection{Soft and collinear contributions}
\label{sec:sc-contributions}

The NNLL corrections associated with the soft-collinear region are
$\delta \mathcal{F}_{\rm sc}$ and $\delta \mathcal{F}_{\rm correl}$. 
Both are expressed in terms of the
soft-collinear limit $\Vsc{\{k_i\}}/v$, defined in Eq.~\eqref{eq:Vsc}, which describes an
ensemble of soft and collinear emissions that are widely separated in
angle.

The correction $\delta\mathcal{F}_{\rm sc}$ corresponds to the addition
of a single soft-collinear emission $k$, parametrised by the standard
variables $\zeta=\Vsc{k}/v$, the emitting leg $\ell$, and its azimuthal
angle $\phi^{(\ell)}$. Physically, this correction accounts for running
coupling corrections to real emissions in the physical scheme and for
the correct rapidity boundary for a single soft-collinear emission. 
\ABquestion{I am unsure what the `the correct rapidity boundary' is?}

The correlated correction $\delta\mathcal{F}_{\rm correl}$ accounts for
the contribution of correlated double emissions. In this case, the additional
soft-collinear emission undergoes a splitting into two partons,
$k_a$ and $k_b$, whose kinematics are described by the variables
introduced in Eq.~\eqref{eq:dFcorrel-general}. This correction therefore 
captures the effect of correlated splittings into either a $q\bar q$ pair or a gluon pair.

The general expression for $\delta\mathcal{F}_{\rm sc}$ is
\begin{multline}
	\label{eq:dFsc-general}
	\delta\mathcal{F}_{\mathrm{sc}}(\lambda)=\frac{\pi}{\alpha_s(Q)}\int_0^\infty\frac{d\zeta}{\zeta}\int_0^{2\pi}\frac{d\phi}{2\pi}\sum_{\ell=1}^3\left[\Delta R'_{\ell,\mathrm{NNLL}}+R''_{\ell,\mathrm{NNLL}}\left(\ln\frac{d_\ell g_\ell(\phi)}{\zeta}-b_\ell\ln\frac{2E_\ell}{Q}\right)\right]\times 
	 \\ \times \int \dZ  \left[\Theta\left(1-\frac{\Vsc{k,\{k_i\}}}{v}\right)
       -\Theta\left(1-\frac{\Vsc{\{k_i\}}}{v}\right)\Theta(1-\zeta)
     \right]\,.
\end{multline}

Using the one-jettiness parameters in Eq.~\eqref{eq:tau1-1em-props}
in the additive result, Eq.~\eqref{eq:dFhc-general-add}, gives
\begin{multline}
\delta\mathcal{F}_{\mathrm{sc}}(\lambda)=-\FNLL(\lambda)\frac{\pi}{\alpha_s(Q)}\sum_{\ell=1}^3\left[\left(\Delta R'_{\ell,\mathrm{NNLL}}+ R''_{\ell,\mathrm{NNLL}}\ln\frac{Q^2}{M Q_{\ell}}\right)\left(\psi^{(0)}\left(1+\RpNLL\right)+\gamma_E\right)\right. \times\\
 \left.+\frac{R''_{\ell,\mathrm{NNLL}}}{2}\left(\left(\psi^{(0)}(1+\RpNLL)+\gamma_E\right)^2-\psi^{(1)}(1+\RpNLL)+\frac{\pi^2}{6}\right)\right]\,,
\end{multline}
where $\RpNLL$ and $R''_{\ell,\mathrm{NNLL}}$ are given in
Eqs.~\eqref{eq:RpNLL} and~\eqref{eq:R2NNLL}, respectively, while
\begin{equation}
\label{eq:DRNNLL}
\Delta R'_{\ell,\rm NNLL}(v) \equiv C_\ell\left(-\alpha_s(Q)\beta_0\,\frac{dg_2^{(\ell)}(\lambda)}{d\lambda}\right)\,.
\end{equation}

The general expression for the correlated correction is
\begin{equation}
\begin{aligned}
\label{eq:dFcorrel-general}
\delta&\mathcal{F}_{\mathrm{correl}}(\lambda)=\sum_{\ell=1}^3\frac{\lambda R''_{\ell,\mathrm{NNLL}}}{2a\beta_0\alpha_s(Q)}
\int_0^\infty\frac{d\zeta}{\zeta}\int_0^{2\pi}\frac{d\phi^{(\ell)}}{2\pi}\int_0^\infty
\frac{d\mu^2}{\mu^2(1+\mu^2)}\int_0^1dz\int_0^{2\pi}\frac{d\phi}{2\pi}\frac{1}{2!}\mathcal{A}^2(z,\mu,\phi)\times\\
&\int\dZ\left[ \Theta\left(1-\frac{V_{\mathrm{sc}}\left(\{\tilde{p}\},k_a,k_b,\{k_i\}\right)}{v}\right) - 
\lim_{\mu^2\to0}\Theta\left(1-\frac{V_{\mathrm{sc}}\left(\{\tilde{p}\},k_a+k_b,\{k_i\}\right)}{v}\right) \right]\,,
\end{aligned}
\end{equation}
where $k=k_a+k_b$, $\mu^2=m^2/k_t^2$ is the rescaled invariant mass of the
splitting, and $\mathcal{A}^2$ is the double-emission function defined in
Ref.~\cite{Arpino:2019ozn}. 

Since one-jettiness is formally equivalent to thrust in
$e^+e^-$ annihilation, both observables being additive with
$a=b_\ell=1$ and $g_\ell(\phi)=1$, we can directly use the angularity
result of Ref.~\cite{Banfi:2018mcq} evaluated at $x=0$:
\ABquestion{What is confusing me is - when do we use 
the work in Ref.~\cite{Banfi:2018mcq} evaluated at $x=0$
and when do we not? The wide-angle correction is not the same 
as in this reference for example. Also is this the correct reasoning
why we can use it: ``both observables being additive with
$a=b_\ell=1$ and $g_\ell(\phi)=1$''?}
\begin{equation}
 	f^{(\ell)}_{\rm correl}\left(z,\mu,\phi,\phi^{(\ell)}\right)=\sqrt{1+\mu^2}\,,
	\label{eq: f-correl-tau_1}
 \end{equation}
which reduces Eq.~\eqref{eq:dFcorrel-general-add} to
\begin{equation}
\label{eq:dFcorrel-tau_1}
	\delta  \mathcal{F}_{\rm correl}(\lambda)=\frac{\pi}{2\alpha_s(Q)}\lambda\,R'' \zeta_2 \FNLL(\lambda)\,.
\end{equation}

\subsection {Hard and collinear contributions}
\label{sec:hc-contributions}

We next consider the NNLL corrections originating from the
hard-collinear region. These corrections arise because a hard-collinear
emission modifies the relation between the underlying Born kinematics
and the accompanying ensemble of soft-collinear emissions. There are
three such contributions:
\begin{itemize}
  \item $\delta\mathcal{F}_{\rm hc}$, which accounts for the emission of
  a hard-collinear parton at the level of the squared matrix element;
  \item $\delta\mathcal{F}_{\rm rec}$, which describes the kinematic
  recoil induced by a hard-collinear emission against the
  soft-collinear ensemble;
  \item $\Delta\mathcal{F}_{\rm rec}$, which accounts for spin
  correlations between a hard-collinear gluon and the event plane.
\end{itemize}
The hard-collinear correction $\delta\mathcal{F}_{\rm hc}$ depends on the
same soft-collinear limit appearing in $\delta\mathcal{F}_{\rm sc}$,
Eq.~\eqref{eq:dFsc-general}. The recoil corrections additionally require
the hard-collinear limit of the observable,
\begin{equation}
  \label{eq:Vhc}
  \frac{V_{\mathrm{hc}}(\{\tilde{p}\}, k_{\mathrm{hc}}, k_1,\dots,k_n)}{v}\equiv \lim_{v\to 0}\frac{\Vfull{k_{\mathrm{hc}}, k_1,\dots,k_n}}{v}\,,
\end{equation}
where $k_{\mathrm{hc}}$ denotes a single hard-collinear emission,
while $k_1,\ldots,k_n$ represent an ensemble independently emitted 
 soft and collinear emissions.
 
We first consider the hard-collinear correction,
\begin{multline}
  \label{eq:dFhc-general}
   \delta\mathcal{F}_{\rm hc}=\sum_\ell \frac{\alpha_s\left(Qv^{1/(a+b_\ell)}\right)}{\alpha_s(Q)\left(a+b_\ell\right)}B_\ell \int_0^{\infty}\frac{d\zeta}{\zeta}\int_0^{2\pi}\frac{d\phi}{2\pi}
 \int \dZ  \times\\ \times  \left[\Theta\left(1-\frac{V_{\mathrm{sc}}(\{\tilde{p}\},k,\{k_i\})}{v}\right) -\Theta(1-\zeta)\Theta\left(1-\frac{V_{\mathrm{sc}}(\{\tilde{p}\},\{k_i\})}{v}\right)\right]\,.
\end{multline}

For one-jettiness, inserting the parameters of
Eq.~\eqref{eq:tau1-1em-props} into the additive expression in
Eq.~\eqref{eq:dFhc-general-add}, gives
 \begin{equation}
  \label{eq:dFhc-tau_1}
    \delta \mathcal{F}_{\rm hc}(\lambda) = \frac{\alpha_s(\sqrt{\tau_1}\,Q)}{2\,\alpha_s(Q)}\frac{d\FNLL}{d\RpNLL} \sum_{\ell}B_\ell\,.
  \end{equation}
For the initial-state legs, the $z$ integration entering
$B_\ell$ (see Eq.~\eqref{eq:Bell-ISR}) cannot be carried out
analytically because of the presence of the PDFs. For the
final-state leg,
$B_3=-\gamma_3^{(0)}$, reproducing the
$e^+e^-$ three-jet result of Ref.~\cite{Arpino:2019ozn}.

The general expression for $\delta\mathcal{F}_{\rm rec}$ is
 \begin{multline}
 \label{eq:dFrec-general}
   \delta\mathcal{F}_{\rm rec}(\lambda)=\sum_\ell \frac{\alpha_s\left(Qv^{1/(a+b_\ell)}\right)}{\alpha_s(Q)\left(a+b_\ell\right)}\int_0^1 dz\,P_\ell(z) \int_0^{\infty}\frac{d\zeta}{\zeta}\int_0^{2\pi}\frac{d\phi}{2\pi}
 \int \dZ  \times\\ \times  \left[\Theta\left(1-\frac{V_{\mathrm{hc}}(\{\tilde{p}\},k_{\mathrm{hc}},\{k_i\})}{v}\right) -\Theta\left(1-\frac{V_{\mathrm{sc}}(\{\tilde{p}\},k,\{k_i\})}{v}\right)\right]\,.
\end{multline}

For one-jettiness, the additive expression,
Eq.~\eqref{eq:dFrec-general-add}, applies. Again, using the results of
Ref.~\cite{Banfi:2018mcq} for angularities at $x=0$, the ratio between the
soft-collinear and hard-collinear measurement functions is
\begin{equation}
\label{eq:tau_1-fhc}
\frac{f_{\mathrm{sc}}^{(\ell)}(z^{(\ell)},\phi^{(\ell)})}{f_{\mathrm{hc}}^{(\ell)}(z^{(\ell)},\phi^{(\ell)})}=
\begin{cases}
z\,, \quad \ell=1,2\,, \\
1/z\,, \quad \ell=3\,.
\end{cases}
\end{equation}
\ABquestion{In this reference we actually get FSR term is $1-z$ as the ratio
so actually don't see how we get this for the ISR or the FSR?
Also why does ISR have $z$ while FSR has $1/z$ cannot see how this makes
sense kinematically speaking?}

Similarly to $\delta\mathcal{F}_{\rm hc}$, the $z$ integration for the
initial-state legs cannot be performed analytically as $P_\ell(z)$ contain PDFs
(Eq.~\eqref{eq:Pell-ISR}). For the final-state leg,
$\ell=3$, the result coincides with the corresponding
$e^+e^-$ three-jet calculation~\cite{Arpino:2019ozn}. Isolating this
contribution, one finds for a final-state quark
\begin{equation}
  \label{eq:dFrec-tau_1q}
  \left.\delta\mathcal{F}_{\rm rec}\right|_{\ell=3}
  =\FNLL\frac{\alpha_s(Q\sqrt{\tau_1})}{2\alpha_s(Q)} C_F\left(\frac{5}{4}-\frac{\pi^2}{3}\right)\,,
\end{equation}
while for a final-state gluon
\begin{equation}
  \label{eq:dFrec-tau_1g}
  \left.\delta\mathcal{F}_{\rm rec}\right|_{\ell=3}
  =\FNLL\frac{\alpha_s(Q\sqrt{\tau_1})}{2\alpha_s(Q)} \left(
 C_A\left(\frac{67}{36}-\frac{\pi^2}{3}\right)
 -T_Rn_f\frac{13}{18} \right)\,.
 \end{equation}
 
The remaining hard-collinear contribution,
$\Delta\mathcal{F}_{\rm rec}$, accounts for spin correlations and is
written as a sum over gluon legs, $\ell_g$,
\ABquestion{Is the sum over legs correct here? Does the $\phi$ and $z$ need to have $\ell_g$ 
superscripts?}
\begin{equation}
\label{eq:DFrec-general}
\begin{aligned}
  \Delta&\mathcal{F}_{\mathrm{rec}}(\lambda) =
\sum_{\ell_g}\frac{\alpha_s(Qv^{1/(a+b_{\ell_g})})}{\alpha_s(Q)(a+b_{\ell_g})}\int_0^\infty \frac{d\zeta}{\zeta}
\int_0^\pi \frac{d\phi}{\pi}\left[2\cos^2\phi^{(\ell_g)} - 1\right]
\int_0^1\!dz\,
\Delta P_g(z^{(\ell_g)})
\\
&\quad \times \int \dZ \left[
\Theta\left(1 - \frac{V_{\mathrm{hc}}(\{\tilde{p}\},k,\{k_i\})}{v}\right)
- \Theta(1-\zeta)\,
\Theta\left(1 - \frac{V_{\mathrm{sc}}(\{\tilde{p}\},\{k_i\})}{v}\right)
\right]\,,
\end{aligned}
\end{equation}

For additive observables, Eq.~\eqref{eq:DFrec-general-add} applies.
Since one-jettiness has no azimuthal dependence,
 the spin-correlation contribution vanishes identically,
\begin{equation}
\Delta\mathcal{F}_{\rm rec}(\lambda)=0\,.
\end{equation}

\ABquestion{I want to explain why this occurs physically - is the following correct:
Spin correlations depend on the orientation of a hard gluon's
polarisation relative to the event plane. An observable can only see
this if it distinguishes radiation in that plane from radiation outside
it. One-jettiness makes no such distinction: it only depends on how
close a hadron is to the beam or to the jet direction, not on its position
around them. With no sensitivity to the event plane, there is nothing
for the polarisation to affect, so $\Delta\mathcal{F}_{\rm rec}$
vanishes identically.}

\subsection{Soft wide angle contributions}
\label{sec:wa-contributions}

The final class of NNLL corrections originates from soft radiation
emitted at wide angles. Two contributions arise in this region.
The first, $\delta\mathcal{F}_{\rm wa}$, accounts for the difference
between the full observable and its soft-collinear approximation for a
single soft wide-angle emission accompanied by an ensemble of
soft-collinear emissions. The second,
$\Delta\mathcal{F}_{\rm wa}$, is associated with resolved soft
wide-angle radiation.

The evaluation of $\delta\mathcal{F}_{\rm wa}$ is naturally expressed in
terms of dipoles $(\ell\bar\ell)$. Accordingly, the momentum
of a single emission $k$ can be decomposed along any pair of light-like 
momenta $(p_{\ell},p_{\bar \ell})$ forming the $(\ell \bar\ell)$ dipole,
\begin{equation}
\label{eq:Sudakov-decomp}
k^{(\ell\bar\ell)}=z^{(\ell)}p_{\ell}+z^{(\bar\ell)}p_{\bar\ell}+\kappa_{\ell\bar\ell}\cos\phi^{(\ell\bar\ell)}n_{\rm{in}}^{(\ell\bar\ell)}+\kappa_{\ell\bar\ell}\sin\phi^{(\ell\bar\ell)}n_{\rm{out}}^{(\ell\bar\ell)}\,,
\end{equation}
where
\begin{equation}
	n_{\rm in}^{(\ell\bar\ell)}=\left(\cot\frac{\theta_{\ell\bar\ell}}{2},\frac{\vec{n}_{\ell}+\vec{n}_{\bar\ell}}{\sin\theta_{\ell\bar\ell}}\right)\,,\quad n_{\rm out}^{(\ell\bar\ell)}=\left(0,\frac{n_{\ell}\times n_{\bar\ell}}{\sin\theta_{\ell\bar\ell}}\right)\,,
	\end{equation}
with $\vec n_{\ell/\bar\ell}=\vec p_{\ell/\bar\ell}/E_{\ell/\bar\ell}$.
The remaining variables entering this decomposition are defined in
Sec.~\ref{subsubsec:Rs}.

The correction $\delta\mathcal{F}_{\rm wa}$ is written in terms of the
soft wide-angle limit of the observable,
\ABquestion{Can we write the limit like this?}
\begin{equation}
  \label{eq:Vwa}
  \frac{V_{\mathrm{wa}}(\{\tilde{p}\}, k^{(\ell\bar{\ell})}, k_1,\dots,k_n)}{v}\equiv \lim_{v\to 0}\frac{\Vfull{k^{(\ell\bar{\ell})}}, k_1,\dots,k_n)}{v}\,,
\end{equation}
where the emission $k^{(\ell\bar\ell)}$ denotes a soft emission at wide angle with
respect to the Born legs.

The corresponding wide-angle correction is
\begin{equation}
\label{eq:dFwa-general}
\begin{aligned}
	&\delta\mathcal{F}_{\mathrm{wa}}(\lambda)=\frac{\alpha_s(v^{1/a}Q)}{a\,\alpha_s(Q)}2\sum_{(\ell\bar\ell)}C_{\ell\bar\ell}\int_0^\infty \frac{d\zeta}{\zeta}
	\int_{-\infty}^\infty d\eta^{(\ell\bar\ell)}\int_0^{2\pi}\frac{d\phi^{(\ell\bar\ell)}}{2\pi}\times\\
	&\times \int \dZ\left[\Theta\left(1-\frac{V_{\mathrm{wa}}\left(\{\tilde{p}\}, k^{(\ell\bar\ell)},\{k_i\}\right)}{v}\right) -\Theta\left(1-\frac{V_{\mathrm{sc}}\left(\{\tilde{p}\}, k^{(\ell\bar\ell)},\{k_i\}\right)}{v}\right)\right]\,,
\end{aligned}
\end{equation}
where the sum runs over all the dipoles $(\ell\bar\ell)$ that build up the Born event. 

For one-jettiness, the additive expression, Eq.~\eqref{eq:dFwa-general-add}, applies. For this observable, we find:
\begin{equation}
\begin{aligned}
	&f_{\rm sc}^{(\ell\bar\ell)}(\eta^{(\ell\bar\ell)},\phi^{(\ell\bar\ell)})=\frac{Q_{\ell\bar\ell}^2}{MQ_{\ell}}e^{-\eta^{(\ell\bar\ell)}}\Theta\left(\eta^{(\ell\bar\ell)}\right)+\frac{Q_{\ell\bar\ell}^2}{MQ_{\bar\ell}}e^{\eta^{(\ell\bar\ell)}}\Theta\left(-\eta^{(\ell\bar\ell)}\right)\,,\\
	&f_{\rm wa}^{(\ell\bar\ell)}(\eta^{(\ell\bar\ell)},\phi^{(\ell\bar\ell)}) =
\min\left\{
\frac{Q_{\ell\bar\ell}^2}{MQ_1}\,e^{-\eta^{(\ell\bar\ell)}},\;
\frac{Q_{\ell\bar\ell}^2}{MQ_2}\,e^{+\eta^{(\ell\bar\ell)}},\;
\frac{Q_{\ell\bar\ell}}{M}\!\left(\cosh(\eta^{(\ell\bar\ell)}-\eta_J)
-\cos(\phi^{(\ell\bar\ell)}-\phi_J)\right)
\right\}\,.
\end{aligned}
\end{equation}
Substituting these expressions into the additive formula yields
\begin{equation}
    \delta\mathcal{F}_{\mathrm{wa}}(\lambda) =
    \mathcal{F}_{\rm NLL}(\lambda)\,
    \frac{\alpha_s(\tau_1 Q)}{\alpha_s(Q)}
    \sum_{(\ell\bar\ell)}C_{\ell\bar\ell}
    \left[\ln^2\frac{Q_{\bar\ell}}{Q_\ell}
    - 2\!\left(I_{\ell\bar\ell,\hat\ell}
    + I_{\bar\ell\ell,\hat\ell}\right)\right]\,,
\end{equation}
where $I_{\ell\bar\ell,\hat\ell}$ are defined in Eq.~(51) of
Ref.~\cite{Alioli:2023rxx} (see also Ref.~\cite{Jouttenus:2011wh}),
and $\hat\ell$ denotes the leg which is neither  $\ell$ or $\bar\ell$.

Unlike $\delta\mathcal{F}_{\rm wa}$, the correction
$\Delta\mathcal{F}_{\rm wa}$ does not involve the wide-angle limit of the
observable, but instead takes the same limits as $\delta\mathcal{F}_{\rm sc}$. The general expression is
%Nevertheless
%it is classified as a soft wide-angle contribution because of the logarithm
%$\ln(Q_{\ell\bar\ell}/Q)$, which originates from soft radiation emitted
%at large angles between the legs $\ell$ and $\bar\ell$. In other words,
%the term ``wide-angle'' refers to the logarithm rather
%than to the kinematic approximation used for the observable.
\begin{multline}
\label{eq:DFwa-general}
	\Delta\mathcal{F}_{\mathrm{wa}}=2\sum_{(\ell\bar\ell)}C_{\ell\bar\ell}\frac{\alpha_s(v^{1/a}Q)}{a\,\alpha_s(Q)}\ln\frac{Q_{\ell\bar\ell}}{Q}\int_0^\infty \frac{d\zeta}{\zeta}
	\int_0^{2\pi}\frac{d\phi}{2\pi}\int \dZ\times \\
	\times \left[\Theta\left(1-\frac{V_{\mathrm{sc}}\left(\{\tilde{p}\}, k,\{k_i\}\right)}{v}\right) -\Theta\left(1-\frac{V_{\mathrm{sc}}\left(\{\tilde{p}\}, \{k_i\}\right)}{v}\right)\left(1-\zeta\right)\right]\,.
\end{multline}

For additive observables, Eq.~\eqref{eq:DFwa-general} simplifies to the
formula of Eq.~\eqref{eq:DFwa-general-add}. Applying this result
to one-jettiness gives
 \begin{equation}
  \label{eq:DFwa-tau_one}
    \Delta \mathcal{F}_{\rm wa}(\lambda) =  \frac{\alpha_s(\tau_1\,Q)}{\alpha_s(Q)}\frac{d\FNLL}{d\RpNLL}2\sum_{(\ell\bar\ell)} C_{\ell\bar\ell}\ln\left(\frac{Q_{\ell\bar\ell}}{Q}\right)\,. 
  \end{equation}
  
This completes the evaluation of all NNLL corrections entering the
resummation of one-jettiness in $Z$+jet production. In the following sections these results are validated against two
independent calculations: the $N$-jettiness subtraction framework of
Ref.~\cite{Gaunt:2015pea} in Sec.~\ref{subsec:N-jet-SCET-val}, and the
Lund-tree shape resummation of Ref.~\cite{vanBeekveld:2025zjh} in
Sec.~\ref{sec:LTS-val}.


 \section{Comparison with $N$-jettiness}
 \label{subsec:N-jet-SCET-val}
 
 Having derived the NNLL cumulant for hadronic collisions, we now validate
our resummation formula, paying particular attention to the treatment of
initial-state radiation. This is achieved by comparing the logarithmic
expansion of the \textsc{ARES} cumulant with the corresponding results
obtained in the SCET-based $N$-jettiness subtraction framework of
Ref.~\cite{Gaunt:2015pea}.

The comparison is performed order by order in $\alpha_s$ and $L$. Expanding the master formula,
$\Sigma_{ij}$, in Eq.~\eqref{eq:Sigma-NNLL}, we write~\cite{Banfi:2004yd}
 \begin{equation}
\label{eq:sig-coeff}
	\Sigma_{ij}(v)=\sum_{n=1}^{\infty}\sum_{m=0}^{n+1} h_{nm}\bar{\alpha}^n_sL^m\,,
\end{equation}
where $\bar{\alpha}_s\equiv\alpha_s(Q)/(2\pi)$. The coefficients
$h_{nm}$ are obtained by expanding the NNLL radiator and collecting
terms order by order in the cumulant. Their explicit expressions for
one-jettiness are given in Appendix~\ref{ap:ARES-coeff}.

These expansion coefficients are compared with the subtraction coefficients,
$\mathcal{C}_n^{(m)}$, appearing in the logarithmic expansion of the
SCET-based $N$-jettiness subtraction formula of
Ref.~\cite{Gaunt:2015pea}. Agreement between the two calculations
provides an independent check of the NNLL resummation derived in this
chapter, and in particular of the treatment of initial-state radiation.

The subtraction coefficients,
$\mathcal{C}_n^{(m)}(\{p\},M,\mu_R)$, are dimensionless functions of the
Born momenta, the invariant mass of the $Z$ boson, and the
renormalisation scale. They are constructed from the hard, beam, jet,
and soft function coefficients, collectively referred to as the subtraction
ingredients.

To facilitate the comparison, we adapt the notation of
Ref.~\cite{Gaunt:2015pea} to the conventions used throughout this
thesis. The incoming partons on legs $\ell=1$ and $\ell=2$ are denoted
by $\kappa_1$ and $\kappa_2$, with flavours $i$ and $j$,
respectively, while the outgoing parton on leg $\ell=3$ is labelled by
$\kappa_3$. The beam-function coefficients
$B_{\kappa_1,n}^{(m)}(x_1,\mu_R,\lambda_F,\lambda_1)$ and
$B_{\kappa_2,n}^{(m)}(x_2,\mu_R,\lambda_F,\lambda_2)$ describe
collinear radiation from the incoming partons with momentum fractions
$x_1$ and $x_2$. The jet-function coefficient
$J_{\kappa_3,n}^{(m)}(\lambda_3)$ describes collinear radiation from the
outgoing parton, while the soft-function coefficient
$\widehat{S}_{\kappa,n}^{(m)}(\{\hat{p}_\ell\},\lambda_s)$, with
$\kappa\equiv\{\kappa_1,\kappa_2,\kappa_3\}$, describes soft radiation
between the coloured legs and depends on the normalised momenta
$\hat{p}_\ell=p_\ell/Q_\ell$. \ABquestion{Should this be $\bar{p}_i?$}

The dimensionless variables entering the subtraction ingredients are
\begin{equation}
    \lambda_3 = \frac{Q_3 M}{\mu_R^2}\,, \qquad
    \lambda_s = \frac{M}{\mu_R}\,, \qquad
    \lambda_1 = \frac{Q_1 M}{\mu_R}\,, \qquad
    \lambda_2 = \frac{Q_2 M}{\mu_R}\,, \qquad
    \lambda_F = \frac{\mu_R^2}{\mu_F^2}\,,
\end{equation}
while their explicit expressions are collected in
Appendix~\ref{ap:subtract-ingred}.

For later convenience, we also adopt the notation of
Ref.~\cite{Gaunt:2015pea} and expand each subtraction ingredient in
logarithms of its corresponding $\lambda$ variable. For example, the
jet function is written as
\begin{subequations}
\label{eq:subt-mn-exp}
\begin{align}
    &J_{\kappa_3,n}^{(m)}(\lambda_3) = J_{\kappa_3,n}^{(m)}
    + \sum_{r=1}^{2m-1-n}\frac{(n+r)!}{n!\,r!}\,J_{\kappa_3,n+r}^{(m)}\,\ln^r\!\lambda_3\,, \\
    &J_{\kappa_3,-1}^{(m)}(\lambda_3) =J_{\kappa_3,-1}^{(m)}+\sum_{n=0}^{2m-1}J_{\kappa_3,n}^{(m)}\,\frac{\ln^{n+1}\!\lambda_3}{n+1}\,, 
\end{align}
\end{subequations}
with analogous expansions holding for the beam and soft functions.

To compare the SCET subtraction coefficients with the \textsc{ARES}
expansion coefficients, we first establish their relation. The
\textsc{ARES} quantities are obtained from the expansion of the
integrated cumulant, whereas those entering the subtraction formula of
Ref.~\cite{Gaunt:2015pea} are defined for the differential distribution.
Accounting also for the different normalisations of the strong coupling,
$2\pi$ versus $4\pi$, gives
\begin{equation}
\label{eq:coeff-convent}
h_{mn}=\frac{\mathcal{C}_{n-1}^{(m)}}{n}\,\frac{(-1)^n}{2m}\,.
\end{equation}
For $\mathcal{C}_{-1}^{(m)}$, the denominator is understood to
correspond to $n=1$.

For brevity, we suppress the arguments of the subtraction coefficients
throughout the remainder of this section. We also partially suppress the
arguments of the beam functions
$B_{\kappa_1,n}^{(m)}(x_1,\lambda_1)$ and
$B_{\kappa_2,n}^{(m)}(x_2,\lambda_2)$, as well as the soft functions
$\widehat{S}_{\kappa,n}^{(m)}(\lambda_s)$.


\subsection{Leading-logarithmic  coefficients}
\label{subsec:LL-coeff}

At LL accuracy, the first subtraction coefficient is
\begin{equation}
\label{eq:C11}
    \mathcal{C}_1^{(1)}=\frac{B_{\kappa_1,1}^{(1)}(x_1,\lambda_1)}{f_i(x_1,\mu_F)}+\frac{B_{\kappa_2,1}^{(1)}(x_2,\lambda_2)}{f_j(x_2,\mu_F)}+
    J_{\kappa_3,1}^{(1)}(\lambda_3)
    +\widehat{S}^{(1)}_{\kappa,1}\left(\lambda_s\right)\,.
\end{equation}
Using the explicit expressions for the individual ingredients given in
Eq.~\eqref{eq:C11-ingred}, this reduces to
\begin{equation}
    \mathcal{C}^{(1)}_1=-4\sum_{\ell}C_{\ell}\,.
\end{equation}
Applying the relation between the SCET subtraction coefficients and the
\textsc{ARES} expansion coefficients in Eq.~\eqref{eq:coeff-convent}, we
obtain
\begin{equation}
    h_{12}=-\sum_{\ell}C_{\ell}\,,
\end{equation}
which agrees with the corresponding \textsc{ARES} expansion coefficient
given in Eq.~\eqref{eq:h12}.

The second subtraction coefficient contributing at LL accuracy is
\begin{equation}
\begin{aligned}
	&\mathcal{C}^{(2)}_3=\frac{B_{\kappa_1,3}^{(2)}(x_1,\lambda_1)}{f_i(x_1,\mu_F)}+\frac{B_{\kappa_2,3}^{(2)}(x_2,\lambda_2)}{f_j(x_2,\mu_F)}+
    J_{\kappa_3,3}^{(2)}(\lambda_3)
    +\widehat{S}^{(2)}_{\kappa,3}\left(\lambda_s\right)\\[0.5ex]
    &\qquad +\left(\frac{B^{(1)}_{\kappa_1,1}(x_1,\lambda_1)}{f_i(x_1,\mu_F)}+\frac{B^{(1)}_{\kappa_2,1}(x_2,\lambda_2)}{f_j(x_2,\mu_F)}\right)\left(J_{\kappa_1,1}^{(1)}(\lambda_3)+\widehat{S}_{\kappa,1}^{(1)}(\lambda_s)\right)+\frac{B^{(1)}_{\kappa_1,1}(x_1,\lambda_1)}{f_i(x_1,\mu_F)}\frac{B^{(1)}_{\kappa_2,1}(x_2,\lambda_2)}{f_j(x_2,\mu_F)}\,.
 \end{aligned}   
\end{equation}
Substituting the explicit expressions from
Eqs.~\eqref{eq:C11-ingred} and \eqref{eq:C23-ingred}, this reduces to
\begin{equation}
    \mathcal{C}^{(2)}_3=8\left(\sum_{\ell}C_{\ell}\right)^2\,.
\end{equation}
Using Eq.~\eqref{eq:coeff-convent}, we then obtain
\begin{equation}
    h_{24}=\frac{1}{2}\left(\sum_{\ell}C_{\ell}\right)^2\,,
\end{equation}
in agreement with the corresponding \textsc{ARES} expansion coefficient
given in Eq.~\eqref{eq:h24}.

\subsection{Next-to-leading-logarithmic coefficients}
\label{subsec:NLL-coeff}

The first subtraction coefficient contributing at NLL accuracy is
\begin{equation}
\label{eq:C10}
    \mathcal{C}_0^{(1)}=\frac{B_{\kappa_1,0}^{(1)}(x_1,\lambda_1)}{f_i(x_1,\mu_F)}+\frac{B_{\kappa_2,0}^{(1)}(x_2,\lambda_2)}{f_j(x_2,\mu_F)}+
    J_{\kappa_3,0}^{(1)}(\lambda_3)
    +\widehat{S}^{(1)}_{\kappa,0}\left(\lambda_s\right)\,.
\end{equation}
Using the explicit expressions for the subtraction ingredients given in
Eqs.~\eqref{eq:C11-ingred} and~\eqref{eq:C10-ingred}, this becomes
\begin{equation}
\label{eq:C10-init}
    \begin{aligned}
    \mathcal{C}^{(1)}_0={}&-2\sum_{\ell} \gamma_{\ell}^{(0)}
    +\sum_{k}\int\frac{dz}{z}2P_{ik}(z)\frac{f_k(x_1/z,\mu_F)}{f_i(x_1,\mu_F)}
    +\sum_{k}\int\frac{dz}{z}2P_{jk}(z)\frac{f_k(x_2/z,\mu_F)}{f_j(x_2,\mu_F)}\\
    &+4C_1\ln\frac{Q_1 M}{\mu_R^2}+4C_2\ln\frac{Q_2 M}{\mu_R^2}
    +4C_3\ln\frac{Q_3 M}{\mu_R^2} \\
    &-4\left[(C_1-C_2-C_3)\ln\frac{Q_{12}^2}{Q_1 Q_2}
    +(C_2-C_1-C_3)\ln\frac{Q_{13}^2}{Q_1 Q_3}
    +(C_3-C_1-C_2)\ln\frac{Q_{12}^2}{Q_1 Q_2}\right]\\
    &-8\ln\frac{M}{\mu_R}\sum_{\ell}C_{\ell}\,,
    \end{aligned}
\end{equation}
which can be simplified to
\begin{equation}
\label{eq:C01-simplified}
    \begin{aligned}
    \mathcal{C}^{(1)}_0={}&-2\sum_{\ell} \gamma_{\ell}^{(0)}
    +\sum_{k}\int\frac{dz}{z}2P_{ik}(z)\frac{f_k(x_1/z,\mu_F)}{f_i(x_1,\mu_F)}
    +\sum_{k}\int\frac{dz}{z}2P_{jk}(z)\frac{f_k(x_2/z,\mu_F)}{f_j(x_2,\mu_F)}\\
    &+4C_1\ln\frac{Q_{12}^2 Q_{13}^2}{Q_{23}^2 MQ_1}
    +4C_2\ln\frac{Q_{12}^2 Q_{23}^2}{Q_{13}^2 MQ_2}
    +4C_3\ln\frac{Q_{13}^2 Q_{23}^2}{Q_{12}^2 MQ_3}\,.
    \end{aligned}
\end{equation}
Applying Eq.~\eqref{eq:coeff-convent}, we obtain the corresponding
\textsc{ARES} expansion coefficient
\begin{equation}
\label{eq:SCET-h11}
    \begin{aligned}
    h_{11}=&\sum_{\ell}\gamma^{(0)}_{\ell}
    -\sum_{k}\int\frac{dz}{z}P_{ik}(z)\frac{f_k(x_1/z,\mu_F)}{f_i(x_1,\mu_F)}
    -\sum_{k}\int\frac{dz}{z}P_{jk}(z)\frac{f_k(x_2/z,\mu_F)}{f_j(x_2,\mu_F)}\\
    &-2C_1\ln\frac{Q_{12}^2 Q_{13}^2}{Q_{23}^2 MQ_1}
    -2C_2\ln\frac{Q_{12}^2 Q_{23}^2}{Q_{13}^2 MQ_2}
    -2C_3\ln\frac{Q_{13}^2 Q_{23}^2}{Q_{12}^2 MQ_3}\,,
    \end{aligned}
\end{equation}
which agrees exactly with the \textsc{ARES} expansion coefficient given
in Eq.~\eqref{eq:h11}.

The remaining subtraction coefficient entering at NLL accuracy is
\begin{equation}
\begin{aligned}
&\mathcal{C}^{(2)}_2=\frac{B_{\kappa_1,2}^{(2)}(x_1,\lambda_1)}{f_i(x_1,\mu_F)}+\frac{B_{\kappa_2,2}^{(2)}(x_2,\lambda_2)}{f_j(x_2,\mu_F)}+
    J_{\kappa_3,2}^{(2)}(\lambda_3)
    +\widehat{S}^{(2)}_{\kappa,2}\left(\lambda_s\right)\\[0.5ex]
    &\qquad +\frac{3}{2}\left[\frac{B_{\kappa_1,0}^{(1)}(x_1,\lambda_1)}{f_i(x_1,\mu_F)}\left(\frac{B_{\kappa_2,1}^{(1)}(x_2,\lambda_2)}{f_j(x_2,\mu_F)}+J_{\kappa_3,1}^{(1)}(\lambda_3)
    +\widehat{S}^{(1)}_{\kappa,1}\left(\lambda_s\right)\right)\right.\\[0.5ex]
    &\qquad\qquad +\frac{B_{\kappa_2,0}^{(1)}(x_2,\lambda_1)}{f_j(x_2,\mu_F)}\left(\frac{B_{\kappa_1,1}^{(1)}(x_1,\lambda_1)}{f_i(x_1,\mu_F)}+J_{\kappa_3,1}^{(1)}(\lambda_3)
    +\widehat{S}^{(1)}_{\kappa,1}\left(\lambda_s\right)\right)\\[0.5ex]
    &\qquad\qquad +J_{\kappa_3,0}^{(1)}(\lambda_3)\left(\frac{B_{\kappa_1,1}^{(1)}(x_1,\lambda_1)}{f_i(x_1,\mu_F)}+\frac{B_{\kappa_2,1}^{(1)}(x_2,\lambda_2)}{f_j(x_2,\mu_F)}+\widehat{S}^{(1)}_{\kappa,1}\left(\lambda_s\right)\right)\\[0.5ex]
    &\qquad\qquad \left.+\widehat{S}^{(1)}_{\kappa,0}\left(\lambda_s\right)\left(\frac{B_{\kappa_1,1}^{(1)}(x_1,\lambda_1)}{f_i(x_1,\mu_F)}+\frac{B_{\kappa_2,1}^{(1)}(x_2,\lambda_2)}{f_j(x_2,\mu_F)}+J_{\kappa_3,1}^{(1)}(\lambda_3)\right)\right]\,.
 \end{aligned}
\end{equation}

Substituting the explicit expressions for the ingredients from
Eqs.~\eqref{eq:C10-ingred},~\eqref{eq:C11-ingred}, and
\eqref{eq:C22-ingred}, the expression reduces to
\begin{equation}
    \mathcal{C}_{2}^{(2)} = 24\pi\beta_0 \sum_{\ell}C_{\ell}
    +12\left(-\frac{\mathcal{C}^{(1)}_{0}}{2}\right)\sum_{\ell}C_{\ell}\,,
\end{equation}
where $\mathcal{C}^{(1)}_0$ is given in Eq.~\eqref{eq:C01-simplified}.

Applying Eq.~\eqref{eq:coeff-convent}, this gives
\begin{equation}
    h_{23}=-2\pi\beta_{0}\sum_{\ell}C_{\ell}-h_{11}\sum_{\ell}C_{\ell}\,,
\end{equation}
which agrees with the corresponding \textsc{ARES} expansion coefficient
given in Eq.~\eqref{eq:h23}.

\subsection{Next-to-next-to-leading-logarithmic coefficients}
\label{subsec:NNLL-coeff}

The first subtraction coefficient entering at NNLL accuracy is
$\mathcal{C}_{-1}^{(1)}$, defined as
\begin{equation}
    \mathcal{C}_{-1}^{(1)}=H^{(1)}_{\rm SCET}\,+\frac{B_{\kappa_1,-1}^{(1)}(x_1,\lambda_1)}{f_i(x_1,\mu_F)}+\frac{B_{\kappa_2,-1}^{(1)}(x_2,\lambda_2)}{f_j(x_2,\mu_F)}+J_{\kappa_3,-1}^{(1)}(\lambda_c)
    +\widehat{S}^{(1)}_{\kappa,-1}\left(\lambda_s\right)\,,
\end{equation}
where the explicit expressions for the individual subtraction ingredients
are given in Eqs.~\eqref{eq:Cd1-ingred},~\eqref{eq:C10-ingred}, and
~\eqref{eq:C11-ingred}. The hard function is given by 
\begin{equation}
\label{eq:CH-SCET}
	H^{(1)}_{\rm SCET}\equiv \frac{\left(\vec{C}^{\dagger(1)}\vec{C}^{(0)}+\vec{C}^{\dagger(0)}\vec{C}^{(1)}\right)}{|\vec{C}^{(0)}|^2}\,,
\end{equation} 
where $\vec{C}^{(0)}$ and $\vec{C}^{(1)}$ denote the tree-level and
one-loop hard Wilson coefficients. \ABquestion{I want
to explain what these are a bit more, can I say: These are obtained by matching the QCD amplitude
onto the corresponding SCET operator basis. These coefficients encode
the short-distance hard interaction and are evaluated at the hard scale.}

\ABquestion{This is my justification to the examiners for why we are allowed to do this - is this OK?:
The hard function is the only process-dependent ingredient entering the
NNLL subtraction coefficient and therefore requires the corresponding
one-loop virtual corrections. Our aim here is not to determine the
physical hard function for $pp\to Z+\mathrm{jet}$, but rather to validate
the relation between the SCET hard function and the corresponding
\textsc{ARES} hard coefficient. This relation is process independent:
the conversion from $H_{\rm SCET}^{(1)}$ to the \textsc{ARES} coefficient
$\mathcal{H}_1$, obtained by subtracting the finite terms of
Eq.~\eqref{eq:HQ-expanded} (excluding $\mathcal{H}_1$), depends only on
the universal subtraction structure determined by the colour charges and
kinematic invariants of the external coloured partons. Since the
$q\bar q\to Hg$ and $q\bar q\to Zg$ channels possess the same
coloured-parton configuration, the known one-loop virtual corrections for
$pp\to H+\mathrm{jet}$ from Ref.~\cite{Ravindran:2002dc} provide a
non-trivial validation of this comparison. The resulting hard
coefficient is naturally specific to Higgs production and is used here
solely for the purpose of validating the formalism.}

The extraction of the $q\bar q\to Hg$ hard function is presented in
Appendix~\ref{ap:H1-extraction}, with the remaining partonic channels
following analogously.


Using Eq.~\eqref{eq:H1-SCET} together with the subtraction ingredients
listed in Eq.~\eqref{eq:Cd1-ingred}, we obtain for the
$q\bar{q}\to Hg$ channel
\begin{equation}
\begin{aligned}
&\mathcal{C}^{(1)}_{-1} = H_{\rm SCET}^{(1)}+\sum_{k}\int\frac{dz}{z} \frac{f_k(x_1/z,\mu_F)}{f_i(x_1,\mu_F)}\left(2I_{qk}^{(1)}(z)+2P_{qk}(z)\ln\frac{Q_1\,M}{\mu_F^2}\right)\\[0.5ex]
&\qquad+\sum_{k}\int\frac{dz}{z} \frac{f_k(x_2/z,\mu_F)}{f_j(x_2,\mu_F)}\left(2I_{\bar{q}k}^{(1)}(z)+2P_{\bar{q}k}(z)\ln\frac{Q_2\,M}{\mu_F^2}\right)\\[0.5ex]
&\qquad -3C_F\ln\frac{Q_1\,M}{\mu_R^2}+2C_F\ln^2\frac{Q_1\,M}{\mu_R^2}-3C_F\ln\frac{Q_2\,M}{\mu_R^2}+2C_F\ln^2\frac{Q_2\,M}{\mu_R^2} \\[0.5ex]
&\qquad + C_A\left(\frac{4}{3}-\pi^2\right)+\frac{20}{3}\pi\beta_0-4\pi\beta_0\ln\frac{Q_3\,M}{\mu_R^2}+2C_A\ln^2\frac{Q_3\,M}{\mu_R^2}\\[0.5ex]
&\qquad+(C_A-2C_F)\left[\ln^2\frac{Q_{12}^2}{Q_1 Q_2}-\frac{\pi^2}{6}+2(I_{12,3}+I_{21,3})\right] \\[0.5ex]
&\qquad-C_A\left[\ln^2\frac{Q_{13}^2}{Q_1 Q_3}-\frac{\pi^2}{6}+2(I_{13,2}+I_{31,2})\right] \\[0.5ex]
&\qquad-C_A\left[\ln^2\frac{Q_{23}^2}{Q_2 Q_3}-\frac{\pi^2}{6}+2(I_{23,1}+I_{32,1})\right] \\[0.5ex]
&\qquad -4\left[(C_A-2C_F)\ln\frac{Q_{12}^2}{Q_1Q_2}-C_A\ln\frac{Q_{13}^2}{Q_1Q_3}-C_A\ln\frac{Q_{23}^2}{Q_2Q_3}\right]\ln\frac{M}{\mu_R}\\[0.5ex]
&\qquad -4(2C_F+C_A)\ln^2\frac{M}{\mu_R}\,.
\end{aligned}
\end{equation}
%
The corresponding \textsc{ARES} expansion coefficient $h_{10}$ for the
$q\bar q\to Hg$ channel, given in Eq.~\eqref{eq:h10}, is
\begin{equation}
\begin{aligned}
&h_{10}=H_{q\bar{q}}^{(1)}+\sum_{k}\int\frac{dz}{z} \frac{f_k(x_1/z,\mu_F)}{f_i(x_1,\mu_F)}\left(I_{qk}^{(1)}(z)+P_{qk}(z)\ln\frac{Q_1\,M}{\mu_F^2}\right)\\[0.5ex]
&\qquad+\sum_{k}\int\frac{dz}{z} \frac{f_k(x_2/z,\mu_F)}{f_j(x_2,\mu_F)}\left(I_{\bar{q}k}^{(1)}(z)+P_{\bar{q}k}(z)\ln\frac{Q_2\,M}{\mu_F^2}\right)\\[0.5ex]
&\qquad-\frac{3}{2}C_F\ln\frac{Q_1\,M}{Q^2}-\frac{3}{2}C_F\ln\frac{Q_2\,M}{Q^2}-2\pi\beta_0\ln\frac{M\,Q_3}{Q^2}\\[0.5ex]
    &\qquad-\frac{2C_F-C_A}{2}\left(\ln^2\frac{Q_{12}^2}{Q_1\,M}+\ln^2\frac{Q_{12}^2}{Q_2\,M}\right)-\frac{C_A}{2}\left(\ln^2\frac{Q_{13}^2}{Q_1\,M}+\ln^2\frac{Q_{13}^2}{Q_3\,M}\right)\\[0.5ex]
&\qquad-\frac{C_A}{2}\left(\ln^2\frac{Q_{23}^2}{Q_2\,M}+\ln^2\frac{Q_{23}^2}{Q_3\,M}\right)+\frac{2C_F-C_A}{2}\left[\ln^2\frac{Q_2}{Q_1}-2(I_{12,3}+I_{21,3})\right]\\[0.5ex]
&\qquad+\frac{C_A}{2}\left[\ln^2\frac{Q_3}{Q_1}-2(I_{13,2}+I_{31,2})\right]+\frac{C_A}{2}\left[\ln^2\frac{Q_3}{Q_2}-2(I_{23,1}+I_{32,1})\right]\\[0.5ex]
&\qquad +\frac{67C_A-20 T_R\,n_f}{18}-\frac{\pi^2}{3}(C_F+C_A)\,,
\end{aligned}
\end{equation}
where we have used the relation
\begin{equation}
\label{eq:I1-ARES}
I_{ik}^{(1)}(z) =
-\hat{P}_{ik}^{(\varepsilon)}(z)
+\hat{P}_{ik}(z;0)\!\left[
\left(\frac{\ln(1-z)}{1-z}\right)_{\!\!+}(1-z)
-\ln z\right]
+\frac{\pi^2}{6}\,\delta_{ij}\,C_i\,\delta(1-z)\,,
\end{equation}
to express the SCET matching functions $I_{qk}^{(1)}(z)$ and
$I_{\bar{q}k}^{(1)}(z)$ in terms of the corresponding \textsc{ARES}
ingredients. The first two terms in Eq.~\eqref{eq:I1-ARES} correspond to
the hard-collinear contribution $C_{\mathrm{hc},\ell}^{(1)}$ and the
recoil correction $\delta\mathcal{F}_{\rm rec}$ associated with
initial-state radiation, as defined in Eqs.~\eqref{eq:Cij-ISR} and
\eqref{eq:dFrec-ISR}.

The remaining comparison between the SCET subtraction coefficient and the
\textsc{ARES} expansion coefficient involves a lengthy rearrangement of
the logarithmic and finite terms and is therefore not shown explicitly
here. Instead, the symbolic implementation of this comparison is provided
in the \textsc{Mathematica} notebook of Ref.~\cite{peake2026validation},
where it is verified that
\begin{equation}
    \mathcal{C}^{(1)}_{-1}=2h_{10}\,,
\end{equation}
as required by Eq.~\eqref{eq:coeff-convent}.

Finally, we consider the remaining subtraction coefficients contributing
at NNLL accuracy, $\mathcal{C}^{(2)}_1$ and $\mathcal{C}^{(2)}_0$, which
correspond to the \textsc{ARES} expansion coefficients $h_{22}$ and
$h_{21}$, respectively\footnote{Note that $h_{22}$ formally enters at
NLL$'$ accuracy.}. Their explicit expressions are given in
Eqs.~\eqref{eq:C21} and~\eqref{eq:C20}. Since these expressions are
considerably more involved than those encountered at LL and NLL
accuracy, we perform the comparison with the \textsc{ARES} expansion
coefficients symbolically using \textsc{Mathematica}. Details of the
implementation are provided in Appendix~\ref{ap:symb-val}, while the
complete notebook is available in Ref.~\cite{peake2026validation}.

In the next section, we provide an alternative analytic validation of
the treatment of initial-state radiation.

\section{Validation: Lund-tree observables}
\label{sec:LTS-val}

In this section we validate the ISR hard-collinear corrections derived in
Sec.~\ref{sec:NNLL-initial} by comparing the \textsc{ARES} master formula in
Eq.~\eqref{eq:Sigma-NNLL} with the NNLL resummation of Ref.~\cite{vanBeekveld:2025zjh} for a class
of observables known as Lund-tree Shapes (LTS). This provides
a particularly clean test of the ISR treatment, since the hadron--hadron
results of Ref.~\cite{vanBeekveld:2025zjh} are restricted to colour-singlet
production without resolved final-state jets ($N=0$). In this Born-level configuration, 
the only coloured legs are the two incoming partons. Consequently, any
difference between the two formulations is entirely associated with the
treatment of initial-state radiation.

Restricting the \textsc{ARES} master formula in Eq.~\eqref{eq:Sigma-NNLL}
 to the same $N=0$ configuration leaves only the two incoming legs 
 $\ell=1,2$, with the only Born channel being $q\bar q\to Z$. The NNLL
cumulant becomes
\begin{multline}
  \label{eq:Sigma-qqbar}
  \Sigma_{q\bar q}(\{p\},v)= e^{-R_{\rm NNLL}(v)} \frac{f_{q}(x_1,v^{\frac{1}{a+b_1}}Q)}{f_{q}(x_1,\mu_F)} \frac{f_{\bar q}(x_2,v^{\frac{1}{a+b_2}}Q)}{f_{\bar q}(x_2,\mu_F)}
  \times \\ \times
  \left(\mathcal{F}^{(q \bar q)}_{\rm NLL}(\lambda)\left(1+\frac{\alpha_s(Q)}{2\pi}H_{q \bar q}^{(1)}(Q,\mu_R)
      +\frac{\alpha_s(Q v^{\frac{1}{a+b_\ell}})}{2\pi}\sum_{\ell}C_{\mathrm{hc},\ell}^{(1)}\right) +\frac{\alpha_s(Q)}{\pi} \delta\mathcal{F}^{(q\bar q)}_{\rm NNLL}(\lambda)\right)\,,
\end{multline}
where 
\begin{equation}
\label{eq:dFNNLL-qq}
	 \delta\mathcal{F}^{(q\bar q)}_{\rm NNLL} = \delta\mathcal{F}_{\rm sc}+\delta\mathcal{F}_{\rm hc}+\delta\mathcal{F}_{\rm wa}+\delta\mathcal{F}_{\rm correl}+\delta \mathcal{F}_{\mathrm{rec}}\,.
\end{equation}
The terms $\Delta\mathcal{F}_{\rm wa}$ and $\Delta\mathcal{F}_{\rm rec}$
are absent from Eq.~\eqref{eq:dFNNLL}. These corrections were introduced
for three-jet configurations in Ref.~\cite{Arpino:2019ozn} and require the
presence of a final-state coloured leg. They therefore do not contribute
to the present $N=0$ configuration.

We restrict ourselves to the additive LTS observable $S_b^{(0)}$,
for which the NNLL resummation contains a comparable ISR
term.\footnote{The other observable they define, $M_b^{(0)}=\max_{j\in\mathcal{I}_0}v_b^j$
is not an event shape as it is not collinear safe and so will not offer a reasonable comparison.}
This allows a term-by-term
comparison of the ISR contributions. The observable is
\begin{equation} 
\label{eq:Sb0-def} 
S_b^{(0)} = \sum_{j\in\mathcal{I}_0}v_b^j \,, \qquad v_b^{j_i} = \frac{p_{t,i}}{M}e^{-b|y_i|}\,,
\end{equation} 
where $(p_{t,i},y_i)$ are the transverse momentum and rapidity of
jet $i$ with respect to the beam axis, $M$ is the invariant mass of
the colour singlet (taken as the hard scale), and
$\mathcal{I}_0\equiv\{j_1,j_2,\dots,j_m\}$ is the list of jets from
the C/A algorithm of radius $R$ (cf.\ Sec.~\ref{sec:gen-kt} for more details on
this algorithm). \ABquestion{Should I change $p_{t,i}, y_i$ to $k_{t^{(\ell)}}, \eta^{(\ell)}$ and have it as a sum over $\ell$?}

Setting $Q=M$ fixes the single-emission parameters to
\begin{equation}
\label{eq:LTS-1-emsn}
	a=1, \quad b_\ell=b, \quad d_\ell=1, \quad g_\ell(\phi)=1,\quad \ell=1,2\,.
\end{equation}

Using the \textsc{ARES} notation and setting $Q=M$, the NNLL cumulant of
Ref.~\cite{vanBeekveld:2025zjh} can be written as
\begin{equation}
\label{eq:LTS-cum}
\begin{aligned}
&\Sigma^{pp}(\{p\},v)=e^{-R_{\mathrm{NNLL}}(v)}
\frac{f_q(x_1,\mu_F v^{\frac{1}{1+b}})f_{\bar q}(x_{2},\mu_F v^{\frac{1}{1+b}})}
{f_q(x_1,\mu_F )f_{\bar q}(x_{2},\mu_F )}\left\{\mathcal{F}_{\mathrm{NLL}}(v)\left[
1+\frac{\alpha_s(\mu_R)}{2\pi}H_{q\bar q}^{(1)}(Q,\mu_R)+
\right.\right.\\[0.5ex]
&\qquad\qquad +\frac{\alpha_s(\mu_R)}{2\pi}\frac{1+b}{1+b-2\lambda}
\sum_k\left(
\int_{x_1}^1\frac{dz}{z}C_{qk}^{(1), \rm LTS}(z)
\frac{f_k\left(x_1/z,\mu_F v^{\frac{1}{1+b}}\right)}{f_{ q}(x_1,\mu_F v^{\frac{1}{1+b}})}\right.\\[0.5ex]
&\qquad\qquad\left.\left.\left.+\int_{x_2}^1\frac{dz}{z}C_{\bar{q}k}^{(1), \rm LTS}(z)
\frac{f_k\left(x_2/z,\mu_F v^{\frac{1}{1+b}}\right)}{f_{\bar{q}}(x_2,\mu_F v^{\frac{1}{1+b}})}\right)\right]+\frac{\alpha_s(\mu_R)}{\pi}\delta\mathcal{F}^{ pp}_{\mathrm{NNLL}}(v)\right\}+\delta\Sigma_{\mathrm{DGLAP}}^{pp}(v)\,,
\end{aligned}
\end{equation}
where 
\begin{equation}
\label{eq:NNLL-LTS}
	\delta\mathcal{F}^{pp}_{\mathrm{NNLL}}(v)=\delta\mathcal{F}^{pp}_{\mathrm{sc}}(v)+\delta\mathcal{F}^{pp}_{\mathrm{wa}}(v)+\delta\mathcal{F}^{pp}_{\mathrm{hc}}(v)+\delta\mathcal{F}^{pp}_{\mathrm{correl}}(v)+\delta\mathcal{F}^{pp}_{\mathrm{clust}}(v)\,,
\end{equation}
and
\begin{equation}
\label{eq:DGLAP-add}
\begin{aligned}
	&\delta\Sigma_{\mathrm{DGLAP}}^{pp}(v)=e^{-R_{\mathrm{NNLL}}(v)}\frac{\alpha_s(\mu_R)}{\pi}\frac{1}{1+b-2\lambda}\frac{d\FNLL}{d\RpNLL}\\
	&\qquad\qquad\times\sum_k\left(
\int_{x_1}^1\frac{dz}{z}P_{qk}(z)
\frac{f_k\left(x_q/z,\mu_F v^{1/(1+b)}\right)}{f_{ q}(x_{1},\mu_F v^{1/(1+b)})}
+\int_{x_2}^1\frac{dz}{z}P_{\bar{q}k}(z)
\frac{f_k\left(x_{2}/z,\mu_F v^{1/(1+b)}\right)}{f_{\bar{q}}(x_{2},\mu_F v^{1/(1+b)})}
\right)\,,
\end{aligned}
\end{equation}
where $P_{qk}(z)$ and $P_{\bar q k}(z)$ are the full Altarelli--Parisi
splitting functions defined in Eq.~\eqref{eq:full-AP}.

The two cumulants have a similar general structure, with the difference being how they organise the NNLL
corrections. The relation between the different contributions in the two formulations can be summarised as follows.

The radiator $R_{\rm NNLL}$, hard coefficient $H^{(1)}_{q\bar{q}}$,
and NLL function $\mathcal{F}_{\rm NLL}$ are identical in the two
formulations. Furthermore, the NNLL corrections of soft origin,
$\delta\mathcal{F}_{\rm sc}$, $\delta\mathcal{F}_{\rm wa}$, and
$\delta\mathcal{F}_{\rm correl}$, are identical. As discussed in
Sec.~\ref{subsec:soft-fact}, these contributions depend only on the
soft-emission kinematics and colour structure, and are therefore
unaffected by whether the emitting leg is in the initial or final
state. They are universal between ISR and FSR and receive no additional
ISR-dependent contributions. These corrections have already been
validated in previous \textsc{ARES} resummations and are treated
identically in Ref.~\cite{vanBeekveld:2025zjh}, so they do not provide
an independent test of the ISR formalism developed here.

The additional correction $\delta\mathcal{F}^{pp}_{\rm clust}$ in
Eq.~\eqref{eq:NNLL-LTS} originates from the dependence of the observable
$S_b^{(0)}$ on the jet-clustering procedure. This is a genuine feature
of jet-rate observables, and the same type of correction appears in
\textsc{ARES} jet-rate resummations~\cite{Banfi:2016zlc}. Since this
term is again of soft origin it is likewise not part of the present validation.

The non-trivial comparison therefore concerns the hard-collinear
contributions associated with initial-state radiation. In \textsc{ARES}
these are organised into $C_{{\rm hc},\ell}^{(1)}$,
$\delta\mathcal{F}_{\rm hc}$, and $\delta\mathcal{F}_{\rm rec}$,
whereas Ref.~\cite{vanBeekveld:2025zjh} introduces
$C_{qk}^{(1),\rm LTS}$, $\delta\mathcal{F}^{pp}_{\rm hc}$, and
$\delta\Sigma^{pp}_{\rm DGLAP}$. Although the two formulations organise
these terms differently, we show below that they contain the same ISR
hard-collinear corrections.

The coefficient function of Ref.~\cite{vanBeekveld:2025zjh}
is\footnote{In Ref.~\cite{vanBeekveld:2025zjh} the final term is
written as $\ln(M^2/\mu_F^2)\hat{P}_{ij}(z;0)$, which appears to be
a misprint since it introduces an uncancelled collinear singularity --
this has been confirmed with the authors of the paper.}
\begin{multline}
  \label{eq:CkiLTS}
  \sum_j\int_{x_1}^1\frac{dz}{z}\,C_{qk}^{(1),\rm LTS}(z)
  \frac{f_k(x_1/z,\mu_F v^{\frac{1}{1+b}})}
  {f_q(x_1,\mu_F v^{\frac{1}{1+b}})}   = \sum_k\int_{x_1}^1\frac{dz}{z}\,
  \frac{f_k(x_1/z,\mu_F v^{\frac{1}{1+b}})}
  {f_q(x_1,\mu_F v^{\frac{1}{1+b}})}
  \left[-\hat{P}_{qk}^{(\varepsilon)}(z)\right.\\
  \left.+\frac{2b}{1+b}\!\left(\!
  \left(\frac{\ln(1-z)}{1-z}\right)_+
  -\frac{\ln z}{1-z}\right)
  (1-z)\hat{P}_{qk}(z;0)
  +\ln\frac{M^2}{\mu_F^2}P_{qk}(z)\right]\,,
\end{multline}
and analogously for $C_{\bar{q}k}^{(1),\mathrm{LTS}}$.
The first two terms correspond to contributions that are treated
separately in \textsc{ARES}, namely the hard-collinear coefficient
$C_{{\rm hc},\ell}^{(1)}$ and the recoil correction
$\delta\mathcal{F}_{\rm rec}$. To see this, we
specialise Eq.~\eqref{eq:Cell-ISR} to the parameters in
Eq.~\eqref{eq:LTS-1-emsn} with $Q=M$ and $E_\ell=M/2$:
\begin{equation}
\begin{aligned}
  \label{eq:Cij-LTS-comp}
  &\frac{\alpha_s(Mv^{\frac{1}{1+b}})}{2\pi}\sum_{\ell=1}^2
  C_{{\rm hc},\ell}^{(1)} =\\
  &\quad=\sum_k\left[\int_{x_1}^1\frac{dz}{z}
  \frac{f_k(x_1/z,Mv^{\frac{1}{1+b}})}
  {f_q(x_1,Mv^{\frac{1}{1+b}})}
  \left(-\hat{P}_{qk}^{(\varepsilon)}(z)
  +\frac{2b}{1+b}
  \!\left(\frac{\ln(1-z)}{1-z}\right)_+\!(1-z)\hat{P}_{qk}(z;0)
  \right)  +\left\{\substack{x_1\leftrightarrow x_2\\[0.5ex] q\leftrightarrow\bar{q}}\right\}\right]\,.
\end{aligned}
\end{equation}
Using Eq.~\eqref{eq:dFrec-ISR} for additive observables with the
parameters of Eq.~\eqref{eq:LTS-1-emsn},
\begin{equation}
  \label{eq:dFrec-LTS-comp}
  \frac{\alpha_s(M)}{\pi}\delta\mathcal{F}_{\rm rec}(\lambda) =
  \frac{\alpha_s(Mv^{\frac{1}{1+b}})}{2\pi}
  \sum_k\!\left[\int_{x_1}^1\frac{dz}{z}
  \frac{f_k(x_1/z,\mu_F v^{\frac{1}{1+b}})}
  {f_q(x_1,\mu_F v^{\frac{1}{1+b}})}
  \left(-\frac{2b}{1+b}\hat{P}_{qk}(z;0)\ln z\right)
  +\left\{\substack{x_1\leftrightarrow x_2\\[0.5ex] q\leftrightarrow\bar{q}}\right\}\right]\,.
\end{equation}
Combining Eqs.~\eqref{eq:Cij-LTS-comp} and~\eqref{eq:dFrec-LTS-comp}
reproduces the first two terms of Eq.~\eqref{eq:CkiLTS} exactly. The
third term, $\ln(M^2/\mu_F^2)P_{qk}(z)$, arises because the PDFs are 
evaluated at different scales in the two formulations:
\textsc{ARES} uses $Mv^{1/(1+b)}$, while Ref.~\cite{vanBeekveld:2025zjh} uses $\mu_Fv^{1/(1+b)}$. The logarithm $\ln(M^2/\mu_F^2)$ is the 
DGLAP evolution between these two scales and compensates for the
difference between the two. The prefactor $(1+b)/(1+b-2\lambda)$ in Eq.~\eqref{eq:LTS-cum} arises from expressing the running coupling $\alpha_s(Mv^{1/(1+b)})$ in terms of $\lambda=\alpha_s(\mu_R)\beta_0\ln(1/v)$ at one loop. 
The remaining hard-collinear correction in the LTS formulation is
\begin{equation}
\label{eq:dFhc-LTS}
\delta\mathcal{F}^{pp}_{\mathrm{hc}}(v)
=-\frac{d\mathcal{F}_{\rm NLL}}{dR_{\rm NLL}}
\frac{1}{1+b-2\lambda}
\sum_{\ell}\gamma_\ell^{(0)}\,.
\end{equation}
Combining this with $\delta\Sigma^{pp}_{\rm DGLAP}$ in Eq.~\eqref{eq:DGLAP-add}, we find that at NNLL this gives rise to:
\begin{multline}
	\frac{\alpha_s(\mu_R)}{\pi}\delta\mathcal{F}^{pp}_{\mathrm{hc}}(v)+
	\delta\Sigma^{pp}_{\rm DGLAP}\simeq\\
	\frac{\alpha_s(\mu_R)}{\pi(1+b)}\frac{1+b}{1+b-2\lambda}\frac{d\mathcal{F}_{\rm NLL}}{dR_{\rm NLL}}\left[\int_{x_1}^1\frac{dz}{z}\left(P_{qk}(z)-\gamma^{(0)}_1\delta_{qk}\right)
\frac{f_k\left(x_1/z,\mu_F v^{1/(1+b)}\right)}{f_{ q}(x_{1},\mu_F v^{1/(1+b)})}
+ \left\{\substack{x_1\leftrightarrow x_2\\[0.5ex] q\leftrightarrow\bar{q}}\right\}
\right]\,.
\end{multline}

This is equivalent to $\delta\mathcal{F}_{\rm hc}$ in \textsc{ARES}, and can see this explicitly by calculating
Eqs.~\eqref{eq:dFhc-ISR} for additive observables,
\begin{equation}
	\frac{\alpha_s(Q)}{\pi}\delta\mathcal{F}_{\rm hc}=\frac{\alpha_s(Qv^{1/(1+b)})}{\pi(1+b)}\frac{d\FNLL}{d\RpNLL}\sum_k \left[\int_{x_1}^1 \frac{dz}{z} \left(P_{qk}(z)-\gamma_1^{(0)}\delta_{qk}\right) \frac{f_k(x_{1}/z,Qv^{1/(a+b_\ell)})}{f_q(x_1,Qv^{1/(a+b_\ell)})}   +\left\{\substack{x_1\leftrightarrow x_2\\[0.5ex] q\leftrightarrow\bar{q}}\right\}\right]\,.
\end{equation}
This is exactly the combination of the DGLAP term with $\delta\mathcal{F}_{\rm hc}^{pp}$. Consequently, the
additional DGLAP term introduced in
Ref.~\cite{vanBeekveld:2025zjh} does not represent a new physical effect, but rather a different organisation of the hard-collinear corrections. Once all terms are combined, the resummed contributions due to ISR in the two formalisms are the same.

\section{Conclusion}
\label{sec:conclusion}

In this chapter we extended the \textsc{ARES} formalism to NNLL resummation in
hadronic collisions, applying it to $hh\to Z+\text{jet}$. The central result is
that the NNLL master formula takes the same structural form as the
$e^+e^-\to3$-jet case of Ref.~\cite{Arpino:2019ozn}, with the radiator and the
NNLL corrections of soft origin carrying over unchanged. What genuinely differs
is the treatment of initial-state radiation, where a collinear splitting on an
incoming leg convolutes with the parton distribution functions, generating
contributions that have no counterpart in $e^+e^-$ annihilation.

The resulting formalism was validated using the additive observable
one-jettiness, chosen because it has already been studied extensively in
the literature and therefore provides a stringent test of the
resummation framework. The validation was performed at two levels.
First, the LL, NLL, and NNLL expansion coefficients were compared
directly with the $N$-jettiness SCET subtraction framework of
Ref.~\cite{Gaunt:2015pea}, finding exact agreement order by order.
Second, the full resummed cumulant, restricted to colour-singlet
production ($N=0$), was compared with the Lund-Tree Shape resummation of
Ref.~\cite{vanBeekveld:2025zjh}. Although the two approaches organise
their hard-collinear contributions differently, we demonstrated that
they contain identical ISR corrections once all terms are combined.
Together, these comparisons provide a comprehensive validation of the
ISR hard-collinear formalism developed in this chapter.

\ABquestion{This next bit is important as it is justifying why
the validation I did was OK - can I confirm this is accurate.}
There are several natural directions in which this work could be
extended. The validation of $h_{10}$ presented in this chapter
establishes the process-independent relation between the SCET hard
function and the corresponding \textsc{ARES} hard coefficient using the
known $q\bar q\to Hg$ one-loop amplitude of
Ref.~\cite{Ravindran:2002dc}. A natural next step would be to repeat this
construction using the physical $q\bar q\to Zg$ hard function, available
for example through \textsc{MCFM}~\cite{Campbell:1999ah,Campbell:2011bn,Campbell:2019dru}.
 or via analytic
continuation of the $e^+e^-\to q\bar q g$ amplitude. This would provide
the process-specific hard coefficient required for phenomenological
predictions of $Z+\mathrm{jet}$ production and complete the ingredients
needed for matching the NNLL resummation to fixed-order calculations.
Further
interesting observables include the additive observable one minus
transverse thrust and the non-additive thrust minor, both of which are
out-of-plane observables sensitive to the spin-correlation effects
derived here. The formalism could also be extended to N$^3$LL accuracy
and to processes involving a larger number of hard coloured legs. The
latter presents a significant challenge, since for
$\ell\geq4$ the wide-angle corrections become matrices in colour space,
substantially increasing the complexity of the calculation. An
extension to deep inelastic scattering should be comparatively
straightforward, as this involves only a single incoming coloured leg. Similarly, applications
to Higgs- or $W$-boson production would both require a dedicated calculation of the 
corresponding process-dependent hard functions.

The broader motivation for this work is precision phenomenology at the
LHC. A general NNLL resummation for hadronic collisions is
particularly valuable given the wide availability of high-order
fixed-order results, since it supplies the resummed ingredient needed
for reliable predictions across the full kinematic range. It also offers a benchmark
for validating parton-shower algorithms aiming to achieve NNLL
accuracy.  Also necessary for monte-carlo validation. 
The next steps towards phenomenological applications are the
matching of the resummation to fixed-order calculations, the inclusion
of hadronisation corrections, and comparisons with experimental data to
extract a value of the strong coupling. 